\section{Virtual actions along symmetric paths}
\label{sec:spt_virtual_path_algebra}

The interpolation in the proof of Theorem~\ref{thm:spt_classification} first
rephases the two endpoint virtual actions to a common factor system and then forms
their direct sum; necessity rests on the discreteness of the class along a
continuous virtual path \cite[Section~II.F.2]{Schuch2011Phases}.

\begin{lemma}[Unitary rephasing of cohomologous virtual actions]
    \label{thm:spt_unitary_rephasing}
    \lean{TNLean.Algebra.ProjectiveRepresentation.rephase,
        TNLean.Algebra.circleScalar_mem_unitaryGroup,
        TNLean.Algebra.ProjectiveRepresentation.rephase_mem_unitaryGroup,
        TNLean.Algebra.ProjectiveRepresentation.exists_unitary_rephase_of_cohomologous,
        TNLean.Algebra.ProjectiveRepresentation.norm_factorSystem,
        TNLean.Algebra.ProjectiveRepresentation.circlePhaseInclusion_eq,
        TNLean.Algebra.ProjectiveRepresentation.exists_unitary_rephase_with_factorSystem}
    \leanok
    \uses{def:projective_rep, def:cohomologous_to}
    Multiplying each matrix $V_g$ of a projective representation with factor system
    $\omega$ by a scalar $\lambda_g\neq0$ gives a projective representation with
    factor system $\lambda_g\lambda_h\lambda_{gh}^{-1}\omega(g,h)$.  Consequently, let
    $V^0$ be a unitary projective representation whose factor system is the
    pointwise phase $\omega_0/|\omega_0|$.  If $\omega_1$ is cohomologous to
    $\omega_0$, then multiplying each $V_g^0$ by a unit-modulus scalar gives a unitary
    projective representation with factor system $\omega_1/|\omega_1|$.
    If two unitary virtual actions on spaces of positive dimension have factor
    systems $\omega_0$ and $\omega_1$, then $|\omega_0(g,h)|=|\omega_1(g,h)|=1$.
    If these systems are cohomologous, one can rephase the first action
    to have the factor system of the second
    \cite[Section~II.F.2, lines~886--921]{Schuch2011Phases}.
\end{lemma}

\begin{proof}\leanok
    Choose $\lambda$ with
    $\omega_1(g,h)=\lambda_g\lambda_h\lambda_{gh}^{-1}\omega_0(g,h)$, and put
    $\psi_g=\lambda_g/|\lambda_g|$.  The unit-phase map $z\mapsto z/|z|$ is
    multiplicative, so
    \begin{align}
        \frac{\omega_1(g,h)}{|\omega_1(g,h)|}
         & =\psi_g\psi_h\psi_{gh}^{-1}\,\frac{\omega_0(g,h)}{|\omega_0(g,h)|}.
        \notag
    \end{align}
    Hence $\psi_gV_g^0$ has factor system $\omega_1/|\omega_1|$, and multiplication
    by a unit-modulus scalar preserves unitarity.
    For a unitary virtual action,
    $1=\|V_gV_h\|=|\omega(g,h)|\,\|V_{gh}\|=|\omega(g,h)|$.
    Hence its factor system equals its pointwise unit phase.
\end{proof}

\begin{lemma}[Fixed-point path for cohomologous virtual actions]
    \label{thm:spt_cohomologous_fixed_point_path}
    \lean{MPSTensor.sptKron_eq_of_eq_smul,
        MPSTensor.sptFixedPointAction_eq_of_forall_eq_smul,
        MPSTensor.exists_symmetricGappedInteractionPath_of_cohomologous_fixedPoint}
    \leanok
    \uses{thm:spt_unitary_rephasing, thm:spt_direct_sum_fixed_point_path,
        def:spt_fixed_point_tensor}
    Let $D_0,D_1>0$, and let $V^0$ and $V^1$ be unitary projective
    representations with cohomologous factor systems $\omega_0$ and
    $\omega_1$. There is a unitary projective representation $V^{0\prime}$
    with factor system $\omega_1$ whose fixed-point physical action, with
    trivial character, equals that of $V^0$.
    Set $D=D_0+D_1$ and $W_g=V^{0\prime}_{g^{-1}}\oplus V^1_{g^{-1}}$.
    Under the common on-site action $U_g=W_g^T\otimes W_g^{-1}$, the
    normalized bond interactions of Lemma~\ref{thm:spt_direct_sum_fixed_point_path}
    give a symmetric gapped path between their two endpoint interactions on
    the physical space of dimension $D^2$
    \cite[Section~II.F.2]{Schuch2011Phases}. The normalized bond follows
    the corrected second-block range of
    \cite{gap:spc11_spt_interpolation_upper_range}.
\end{lemma}

\begin{proof}\leanok
    Choose $\lambda$ with
    $\omega_1(g,h)=\lambda_g\lambda_h\lambda_{gh}^{-1}\omega_0(g,h)$.
    Lemma~\ref{thm:spt_unitary_rephasing} gives
    $V^{0\prime}_g=(\lambda_g/|\lambda_g|)V^0_g$ with factor system
    $\omega_1$. Scalars cancel in $W^T\otimes W^{-1}$, so the first
    fixed-point physical action is unchanged.
    Lemma~\ref{thm:spt_direct_sum_fixed_point_path} now applies to
    $V^{0\prime}\oplus V^1$.
\end{proof}

\begin{lemma}[Local continuous conjugators from matrix units]
    \label{thm:spt_local_matrix_conjugator}
    \lean{Matrix.matrixUnitConjugator, Matrix.matrixUnitConjugator_eq_smul,
        Matrix.matrixUnitConjugator_intertwines,
        Matrix.exists_continuous_local_matrixConjugator}
    \leanok
    Let $\alpha_t$ be a family of inner actions on $M_D(\C)$,
    with $D>0$, such that $t\mapsto\alpha_t(M)$ is continuous
    for every $M$. Suppose $\alpha_{t_0}=\operatorname{Ad}X_0$
    with $X_0$ invertible. Fix $a\in\{1,\ldots,D\}$ and set
    \begin{align}
        X(t)_{ij}=(\alpha_t(E_{ja})X_0)_{ia}.
        \notag
    \end{align}
    Then $X$ is continuous, $X(t_0)=X_0$, and
    $\alpha_t(M)X(t)=X(t)M$. It is invertible in a neighborhood
    of $t_0$. This is an auxiliary local construction for
    \cite[Section~II.F.2]{Schuch2011Phases}.
\end{lemma}

\begin{proof}\leanok
    If $\alpha_t=\operatorname{Ad}Y$, matrix-unit multiplication gives
    $X(t)=(Y^{-1}X_0)_{aa}Y$. This proves the intertwining identity,
    and gives $X(t_0)=X_0$. Continuity follows entrywise from the
    formula. Its determinant is nonzero near $t_0$.
\end{proof}

\begin{lemma}[Continuous inverse coefficients of an injective tensor]
    \label{thm:spt_continuous_inverse_coefficients}
    \lean{MPSTensor.coefficientInverseMatrix,
        MPSTensor.physicalMatrix_transpose_mul_coefficientInverseMatrix,
        MPSTensor.continuous_coefficientInverseMatrix_family,
        MPSTensor.coefficientInverseMatrix_reconstruct,
        MPSTensor.virtualActionOfTensors, MPSTensor.virtualActionOfTensors_eq_conj,
        MPSTensor.continuous_virtualActionOfTensors_family}
    \leanok
    Let $P$ be the physical matrix of a one-site injective tensor $A$.
    Then $P^\dagger P$ is positive definite, and
    \begin{align}
        C_A=((P^\dagger P)^{-1}P^\dagger)^T,
        \qquad P^T C_A=I.
        \notag
    \end{align}
    Thus $C_A$ expands every virtual matrix in the tensor letters.
    It depends continuously on a continuous injective family.
    Replacing the letters of $A$ by gauge-equivalent letters of $B$
    in this expansion defines their virtual conjugation action,
    also continuously. This gives the inverse of the injective
    tensor map used in \cite[Section~II.F.2]{Schuch2011Phases}.
\end{lemma}

\begin{proof}\leanok
    Injectivity gives full column rank of $P$, so $P^\dagger P$ is
    positive definite. Transpose the left-inverse identity to obtain
    $P^TC_A=I$. Inversion is continuous on invertible matrices.
    If $B^i=YA^iY^{-1}$, the coefficient expansion of $M$ becomes
    $YMY^{-1}$ after replacing the letters.
\end{proof}

\begin{lemma}[Local virtual gauges from an exact tensor symmetry]
    \label{thm:spt_local_virtual_gauge}
    \lean{MPSTensor.exists_continuous_local_gauge_of_gaugeEquiv,
        MPSTensor.exists_continuous_local_virtualGauge_of_isOnSiteSymmetric}
    \leanok

    A continuous pair of pointwise gauge-equivalent one-site injective
    tensor families of fixed positive bond dimension admits locally
    continuous invertible bond gauges near any prescribed initial gauge.
    In particular, if a continuous injective family has exact on-site
    symmetry, then for each group element its virtual symmetry matrix
    can be chosen continuously and invertibly near each parameter.
    No continuous virtual family is assumed.
    This is the local, fixed-dimension step of
    \cite[Section~II.F.2]{Schuch2011Phases}; interval gluing and changes
    of minimal bond dimension remain separate steps, as recorded in
    \cite{gap:spc11_spt_interpolation_upper_range}.
\end{lemma}

\begin{proof}\leanok
    \uses{thm:spt_local_matrix_conjugator,
        thm:spt_continuous_inverse_coefficients}
    Construct the continuous inner action from inverse tensor
    coefficients, then apply the matrix-unit conjugator formula.
    Exact on-site symmetry gives pointwise gauge equivalence by the
    injective fundamental theorem, so the argument applies to each twist.
\end{proof}

\begin{lemma}[Continuous gauges on an interval]
    \label{thm:spt_global_virtual_gauge}
    \lean{MPSTensor.gauge_covariance_of_invertible_intertwiner,
        MPSTensor.exists_units_smul_of_invertible_intertwiners,
        MPSTensor.exists_continuous_gauge_on_unitInterval}
    \leanok

    Let $A_t,B_t$, $t\in[0,1]$, be continuous, pointwise
    gauge-equivalent tensor families of the same positive bond dimension,
    with $A_t$ one-site injective. Given an invertible initial gauge $X_0$,
    there is a continuous family of invertible matrices $X_t$ with that
    initial value and $B_t^iX_t=X_tA_t^i$ for every $t,i$.
    This proves the interval lifting of the local construction in
    \cite[Section~II.F.2]{Schuch2011Phases}.
\end{lemma}

\begin{proof}\leanok
    \uses{thm:spt_local_virtual_gauge}
    The neighborhoods where local gauges are invertible cover the
    interval. Choose a finite subdivision subordinate to this cover.
    At each junction, injectivity implies that the two invertible
    gauges differ by a nonzero scalar. Rescale the next gauge by that
    constant and glue the continuous pieces. The same scalar matching
    gives the prescribed initial gauge.
\end{proof}

\begin{lemma}[A constant character on continuous invariant ground lines]
    \label{thm:spt_ground_line_character}
    \lean{Matrix.mul_rankOne_idempotent_eq_trace_smul,
        Matrix.trace_mul_mem_spectrum_of_rankOne_idempotent_commute,
        Matrix.trace_mul_constant_of_continuous_rankOne_idempotent,
        Matrix.exists_character_of_rankOne_idempotent_commute,
        Matrix.exists_constant_character_of_continuous_rankOne_idempotent}
    \leanok
    Let $P_t$ be a continuous family of rank-one idempotents on a
    finite-dimensional complex space, with parameter in a connected
    space. Suppose each $P_t$ commutes with a fixed unitary representation
    $U$ of a group $G$. There is one unit-modulus character $\chi$ such that
    $U_gP_t=\chi(g)P_t$ for every $g,t$.
    No finiteness assumption on $G$ is needed. Orthogonal projections
    onto unique ground-state lines are a special case, provided their
    continuity has been established. This is an ingredient for the phase
    gauges allowed in \cite[Section~II.C.2 and Appendix~B]{Schuch2011Phases}.
\end{lemma}

\begin{proof}\leanok
    A matrix $X$ commuting with a rank-one idempotent $P$ satisfies
    $XP=\operatorname{tr}(XP)P$. This scalar belongs to the finite
    spectrum of $X$. For fixed $U_g$, its trace against $P_t$ varies
    continuously and hence is constant on the connected parameter space.
    The representation law and $P_t\ne0$ give the character law.
    Unitarity puts each scalar on the unit circle.
\end{proof}

\begin{lemma}[A one-site character from chain lengths two and three]
    \label{thm:spt_character_from_positive_lengths}
    \lean{MPSTensor.mpv_blockTensor_eq_smul_of_mpv_eq_smul_from_two,
        MPSTensor.eq_pow_div_of_mpv_eq_smul_from_two,
        MPSTensor.exists_character_pow_of_mpv_eq_smul_from_two,
        MPSTensor.exists_character_sameMPV_smul_of_mpv_eq_smul_from_two}
    \leanok
    \uses{def:injective}
    Let $B$ be one-site injective with positive bond dimension and suppose
    $\psi_N(A)=c_N\psi_N(B)$ for every $N\ge2$, with $|c_N|=1$.
    Then $c_N=(c_3/c_2)^N$ for every $N\ge2$.
    If the coefficients arise from group characters $c_N(g)$, the ratio
    $\chi(g)=c_3(g)/c_2(g)$ is a unit-modulus character, and
    $c_N(g)=\chi(g)^N$.
    When $A(g)$ and $B$ have the same bond dimension, the tensors
    $\chi(g)B$ and $A(g)$ have equal periodic vectors at all lengths,
    including length one. This gives the one-site phase gauge in
    \cite[Appendix~B]{Schuch2011Phases} from rings of length at least two.
\end{lemma}

\begin{proof}\leanok
    For each $L\ge2$, block into words of length $L$. Every positive
    blocked length corresponds to an original length at least two.
    The injective proportionality theorem gives $c_{NL}=c_L^N$.
    Comparing $2N$ and $3N$ in the two orders gives
    $c_N^2=c_2^N$ and $c_N^3=c_3^N$. Division yields the claimed ratio.
    The character law follows from those of $c_2,c_3$.
    In equal bond dimension, injectivity extends the rephased agreement
    from lengths at least two to every length.
\end{proof}

\begin{lemma}[Removing a physical symmetry character]
    \label{thm:spt_physical_character_twist}
    \lean{MPSTensor.physicalCharacterTwist,
        MPSTensor.physicalCharacterTwist_apply,
        MPSTensor.physicalCharacterUntwist,
        MPSTensor.physicalCharacterUntwist_apply,
        MPSTensor.twistedTensor_physicalCharacterTwist,
        MPSTensor.twistedTensor_physicalCharacterUntwist,
        MPSTensor.isOnSiteSymmetric_physicalCharacterUntwist}
    \leanok
    Multiplying an on-site unitary representation $U_g$ by a unit-modulus
    character $\chi(g)$ gives another unitary representation.
    Tensor twisting then acquires the scalar $\chi(g)$.
    Consequently, if the twisted periodic vectors of $A$ equal those of
    $\chi(g)A$ at every length, replacing $U_g$ by $\chi(g)^{-1}U_g$
    makes the symmetry exact at every length.
    This is the physical phase-gauge change allowed in
    \cite[Section~II.C.2 and Appendix~B]{Schuch2011Phases}.
\end{lemma}

\begin{proof}\leanok
    Scalars commute with the matrix action, and the character law gives
    multiplicativity. Unit modulus preserves unitarity. Pulling the scalar
    out of the physical sum scales the tensor, and hence scales its
    length-$N$ vector by $\chi(g)^N$. The inverse character cancels it.
\end{proof}

\begin{lemma}[Orthogonal original endpoint spaces in the common representation]
    \label{thm:spt_common_physical_endpoint_orthogonality}
    \lean{MPSTensor.commonPhysicalEmbeddingLeft_conjTranspose_mul_right_eq_zero}
    \leanok

    The prescribed common physical embeddings of the two original tensor
    spaces have orthogonal ranges. This holds for arbitrary original
    tensors, without injectivity or symmetry assumptions, and gives
    the orthogonal endpoint summands required by
    \cite[Section~II.C.2]{Schuch2011Phases}.
\end{lemma}

\begin{proof}\leanok
    \uses{thm:spt_common_physical_endpoints}
    The two active coordinate frames occupy disjoint bond blocks.
    Their retained orthogonal complements occupy the two separate
    original physical summands. The mixed adjoint product therefore
    reduces to the zero mixed product of the active frames.
\end{proof}

\begin{lemma}[A gapped fixed-point path retaining the original physical spaces]
    \label{thm:spt_common_physical_fixed_point_path}
    \lean{Matrix.commute_singleKrausMap_of_unitary_intertwiner,
        MPSTensor.isometricInteractionExtension_commute_onSiteTensorPow,
        MPSTensor.normalizedBondInteraction_translate_commute,
        MPSTensor.commonPhysicalNormalizedBondInteraction,
        MPSTensor.commonPhysicalNormalizedBondInteraction_isStarProjection,
        MPSTensor.continuous_commonPhysicalNormalizedBondInteraction,
        MPSTensor.isometricInteractionExtension_ker_le_of_ker_le,
        MPSTensor.commonPhysicalNormalizedBondInteraction_zero_mem_spectrum,
        MPSTensor.commonPhysicalNormalizedBondInteraction_le_parent,
        MPSTensor.commonPhysicalNormalizedBondInteraction_translate_commute,
        MPSTensor.commonPhysicalNormalizedBondInteraction_spectrum_gap_one,
        MPSTensor.commonPhysicalNormalizedBondInteraction_groundSpace,
        MPSTensor.toEuclideanLin_sitewisePhysicalMatrix_mpv_rotatePhysical,
        MPSTensor.commonPhysicalNormalizedBondInteraction_groundSpace_eq_span_mpv,
        MPSTensor.commonPhysicalNormalizedBondInteraction_ker_le_parent,
        MPSTensor.commonPhysicalParentInteraction_commute_onSite,
        MPSTensor.commonPhysicalNormalizedBondInteraction_commute_onSiteTensorPow,
        MPSTensor.commonPhysicalNormalizedBondInteraction_commute_onSite,
        MPSTensor.commonPhysicalNormalizedBondGappedPath,
        MPSTensor.commonPhysicalBondCanonicalParentComparisonPath,
        MPSTensor.commonPhysicalLeftParentComparisonPath,
        MPSTensor.commonPhysicalRightParentComparisonPath,
        MPSTensor.commonPhysicalCanonicalFixedPointGappedPath}
    \leanok

    Let two positive-dimensional unitary virtual representations have the
    same factor system. In the common physical representation consisting
    of their direct-sum adjoint action and the two original physical
    representations, the normalized bond interpolation gives a continuous
    symmetric interaction path. Each periodic Hamiltonian, at every length
    $N\ge2$, has ground energy zero, gap at least one, and ground space equal
    to the periodic MPS line of the included weighted matrix-unit tensor.
    The unused physical states have energy one. Symmetric affine comparisons
    join its endpoints to the canonical two-site parents of the two included
    fixed-point tensors. This realizes the common-space construction in
    \cite[Section~II.F.2]{Schuch2011Phases} for the stated two-site setting.
\end{lemma}

\begin{proof}\leanok
    \uses{thm:spt_common_physical_endpoints,
        thm:spt_spectator_bond_transport,
        thm:spt_normalized_bond_local_symmetry}
    The enlarged local term is
    $I-Q^{\otimes2}(I-h)Q^{\otimes2\dagger}$, where $Q$ includes the active
    physical summand. Resolving each site into the active summand and its
    orthogonal complement shows that the translated projections commute.
    Their sum has spectrum in $\{0\}\cup[1,\infty)$, and its kernel is the
    included old ground line. Physical covariance gives symmetry.
    The ordered affine comparisons preserve the common zero modes and a
    uniform positive gap; concatenate them with the bond interpolation.
\end{proof}

\begin{lemma}[The common-space bond path has an independent exact MPS realization]
    \label{thm:spt_common_physical_fixed_point_ground_path}
    \lean{MPSTensor.normalizedBondInterpolationMatrix_diagonal_ne_zero,
        MPSTensor.weightedMatrixUnitInterpolation_isInjective,
        MPSTensor.mpv_weightedMatrixUnitInterpolation_ne_zero,
        MPSTensor.mpv_rotatePhysical_isometry_ne_zero,
        MPSTensor.commonPhysicalWeightedMatrixUnitInterpolation_injective_representative,
        MPSTensor.exactMPSGroundPath_commonPhysicalNormalizedBondInteraction,
        MPSTensor.commonPhysicalNormalizedBondExactMPSGroundPath}
    \leanok
    \uses{def:spt_exact_mps_gapped_phase}
    For positive endpoint bond dimensions, the included weighted
    matrix-unit interpolation is a continuous exact MPS realization of
    the common-space normalized bond path. Its periodic vectors are
    nonzero at every positive length. At each interior parameter the
    tensor is one-site injective on the full direct-sum bond space; at
    the endpoints it has the corresponding smaller injective
    representative with the same positive-length periodic rays.
    Its finite-chain ground eigenspaces are precisely those periodic
    MPS lines. The ground-state certificate itself requires no
    virtual representation; symmetry comes from the physical path.
    This verifies the exact MPS interpolation in
    \cite[Section~II.F.2]{Schuch2011Phases} in the stated setting.
\end{lemma}

\begin{proof}\leanok
    \uses{thm:spt_common_physical_fixed_point_path}
    Every weighted matrix-unit coefficient is nonzero in the interior,
    so its letters span the full matrix algebra. At an endpoint, the
    surviving corner is the smaller fixed-point tensor. A diagonal word
    has nonzero trace, giving nonvanishing at every positive length;
    an isometric physical inclusion preserves this property.
    The previously computed kernel is exactly the included MPS line,
    with least energy zero.
\end{proof}

\begin{lemma}[Local on-site commutation gives periodic symmetry]
    \label{thm:spt_interaction_hamiltonian_symmetry}
    \lean{MPSTensor.interactionHamiltonian_commute_onSiteTensorPow}
    \leanok
    If a two-site interaction commutes with the tensor square of an
    on-site matrix, then its translation-invariant periodic Hamiltonian
    commutes with the full tensor power at every chain length at least
    two. Unitarity is not needed for this algebraic identity.
\end{lemma}

\begin{proof}\leanok
    Each embedded interaction term has the asserted commutation
    identity. Summing over its translates gives the result.
\end{proof}

\begin{lemma}[Aligning unitary virtual factors]
    \label{thm:spt_unitary_factor_alignment}
    \lean{TNLean.Algebra.ProjectiveRepresentation.norm_factor_eq_one_of_unitary,
        TNLean.Algebra.ProjectiveRepresentation.circlePhaseInclusion_eq_of_unitary,
        TNLean.Algebra.ProjectiveRepresentation.exists_unitary_rephase_of_unitary_cohomologous}
    \leanok
    A unitary projective representation on a nonzero finite-dimensional
    space has a unit-modulus factor system. Two such representations with
    cohomologous factor systems can therefore be rephased by unit-circle
    scalars to have precisely the same factor system, preserving unitarity
    \cite[Section~II.F.2]{Schuch2011Phases}.
\end{lemma}

\begin{proof}\leanok
    Take operator norms in the projective multiplication law. Both unitary
    factors have norm one, so the scalar factor does also. Apply the
    unit-circle representative of the cohomology relation to rephase the
    second virtual representation.
\end{proof}

\begin{lemma}[Physical character changes preserve the same gapped path]
    \label{thm:spt_character_gapped_path}
    \lean{Matrix.finKronecker_const_smul,
        MPSTensor.SymmetricGappedInteractionPath.physicalCharacterTwist,
        MPSTensor.ExactMPSGroundPath.physicalCharacterTwist}
    \leanok
    \uses{def:spt_exact_mps_gapped_phase}
    Multiplying the physical representation by a unit-modulus character
    preserves the same symmetric gapped interaction path and the same
    exact MPS ground realization. Both the Hamiltonians and their ground
    spaces are unchanged. This is the physical phase-gauge freedom of
    \cite[Section~II.C.2 and Appendix~B]{Schuch2011Phases}.
\end{lemma}

\begin{proof}\leanok
    \uses{thm:spt_physical_character_twist}
    The length-$N$ physical action is multiplied by $\chi(g)^N$.
    Scalar multiplication preserves commutation with the Hamiltonian.
    Every other condition concerns the same interactions and tensor family.
\end{proof}

\begin{lemma}[Exact MPS ground lines along a canonical injective path]
    \label{thm:spt_canonical_injective_ground_path}
    \lean{MPSTensor.interactionHamiltonian_parent_groundSpace_eq_span_mpv,
        MPSTensor.exactMPSGroundPath_of_canonicalInjective}
    \leanok
    \uses{def:spt_exact_mps_gapped_phase}
    A continuous one-site injective tensor family of fixed positive bond
    dimension certifies the exact MPS ground lines of any physical path
    whose local interactions are its canonical two-site parents.
    At every periodic length at least two the parent kernel is precisely
    the MPS line, with least energy zero. The periodic vectors are also
    nonzero at length one. No additional virtual data are required by
    this ground-state certificate.
\end{lemma}

\begin{proof}\leanok
    One-site injectivity gives nonvanishing at each positive length.
    The periodic parent-ground-space theorem identifies the kernel
    with the range of the scalar-to-state map, hence the MPS line.
    Positivity gives the least energy zero. Use the given tensor family
    itself as the pointwise injective representative.
\end{proof}

\begin{definition}[Supported isometric bond lifts]
    \label{def:spt_supported_bond_lift}
    \lean{MPSTensor.isometricBondEmbedding, MPSTensor.supportedBondLift}
    \leanok
    Let $V:\C^D\to\C^E$ be an isometry, let $A_\gamma$ be a tensor
    family, and let $T$ be a tensor on $\C^E$. The isometric bond
    inclusion of $A$ has letters $VA^iV^\dagger$. The supported lift
    with prescribed initial tensor $T$ has letters
    \begin{align}
        B_\gamma^i & =VA_\gamma^iV^\dagger
        +T^i-VA_0^iV^\dagger.
        \notag
    \end{align}
\end{definition}

\begin{lemma}[Supported bond lifts preserve the exact ground realization]
    \label{thm:spt_supported_bond_lift}
    \lean{MPSTensor.mpv_isometricBondEmbedding,
        MPSTensor.continuous_isometricBondEmbedding,
        MPSTensor.supportedBondLift_zero,
        MPSTensor.supportedBondLift_intertwine,
        MPSTensor.supportedBondLift_rowSupport,
        MPSTensor.supportedBondLift_samePositiveMpvRay,
        MPSTensor.continuousOn_supportedBondLift,
        MPSTensor.ExactMPSGroundPath.of_supported_isometric_bond_intertwiner,
        MPSTensor.ExactMPSGroundPath.supportedBondLift}
    \leanok
    \uses{def:spt_supported_bond_lift, def:spt_exact_mps_gapped_phase}
    Suppose $T^iV=VA_0^i$ and $VV^\dagger T^i=T^i$ for every letter.
    The supported lift satisfies $B_0=T$, depends continuously on a
    continuous family $A_\gamma$, and has exactly the same periodic
    vectors as $A_\gamma$ at every positive length. Consequently it
    preserves an exact MPS ground realization of the same physical path,
    including its pointwise injective representatives.
\end{lemma}

\begin{proof}\leanok
    The correction annihilates $V$ on its right and is supported on
    the range of $V$ on its left. Thus $B_\gamma^iV=VA_\gamma^i$
    and $VV^\dagger B_\gamma^i=B_\gamma^i$. Multiplying along any
    nonempty word and taking the trace gives the original periodic
    vector. The correction is constant in the parameter, so continuity
    is preserved. All ground-line conditions concern these same vectors.
\end{proof}

\begin{lemma}[Exact ground lines through the canonical fixed-point comparisons]
    \label{thm:spt_common_physical_canonical_ground_path}
    \lean{MPSTensor.interactionHamiltonian_affine_ker_eq_of_order,
        MPSTensor.ExactMPSGroundPath.of_ordered_affine,
        MPSTensor.commonPhysicalBondCanonicalParentComparisonExactMPSGroundPath,
        MPSTensor.commonPhysicalLeftParentComparisonExactMPSGroundPath,
        MPSTensor.commonPhysicalRightParentComparisonExactMPSGroundPath,
        MPSTensor.ExactMPSGroundPath.trans_tensor_zero,
        MPSTensor.ExactMPSGroundPath.trans_tensor_one,
        MPSTensor.commonPhysicalCanonicalFixedPointExactMPSGroundPath,
        MPSTensor.commonPhysicalCanonicalFixedPointExactMPSGroundPath_tensor_zero,
        MPSTensor.commonPhysicalCanonicalFixedPointExactMPSGroundPath_tensor_one}
    \leanok

    The symmetric gapped path between the canonical parents of the two
    included fixed-point tensors has an independent exact MPS ground
    realization. Its initial and final tensors are precisely the included
    weighted endpoint tensors. The affine parent comparisons use the
    constant ground tensor at the corresponding endpoint.
    More generally, an ordered affine comparison between positive
    Hamiltonians with the same zero modes preserves that ground line.
    This certifies the parent comparisons in
    \cite[Section~II.F.2]{Schuch2011Phases}.
\end{lemma}

\begin{proof}\leanok
    \uses{thm:spt_common_physical_fixed_point_ground_path,
        def:spt_exact_ground_path_concatenation}
    If $0\le A\le B$ and $\ker A\subseteq\ker B$, positivity implies
    equality of the two kernels. Every ordered affine combination has
    the same kernel and least energy zero. Keep its MPS tensor constant.
    At the two joins these tensors coincide with the endpoint
    tensors of the bond interpolation. Concatenate the three exact
    ground realizations; the endpoint equations follow from their
    common parameterization.
\end{proof}

\begin{lemma}[Cohomologous unitary factors give a common-space path]
    \label{thm:spt_common_physical_injective_path}
    \lean{MPSTensor.commonPhysicalInjectiveGappedPath,
        MPSTensor.rotatePhysical_covariance_of_circle_rephase,
        MPSTensor.exists_commonPhysicalInjectiveGappedPath_of_cohomologous}
    \leanok

    Let two one-site injective tensors of positive bond dimension have
    unitary physical representations and unitary virtual projective
    representations giving their adjoint tensor covariances. If the
    virtual factor systems are cohomologous, then the canonical
    two-site parents of their prescribed common physical images are
    joined by a symmetric uniformly gapped interaction path.
    The embeddings retain both entire original physical spaces.
    No chosen cochain or additional target intertwiner is an assumption.
    This is the one-site injective, exactly covariant construction from
    \cite[Section~II.F.2]{Schuch2011Phases}. The independent exact MPS
    ground realization is constructed in
    Definition~\ref{def:spt_common_physical_injective_ground_path}.
\end{lemma}

\begin{proof}\leanok
    \uses{thm:spt_common_physical_polar_endpoints,
        thm:spt_common_physical_fixed_point_path,
        thm:spt_unitary_factor_alignment}
    A unit-circle rephasing aligns the factors and cancels out
    of adjoint tensor covariance. The two common physical polar paths
    therefore reach the prescribed included fixed-point tensors.
    Join their canonical parents by the common-space fixed-point path,
    reverse the first polar path, and concatenate.
\end{proof}

\begin{definition}[Concatenating exact MPS ground realizations]
    \label{def:spt_exact_ground_path_concatenation}
    \lean{MPSTensor.ExactMPSGroundPath.trans}
    \leanok
    \uses{def:spt_exact_mps_gapped_phase}
    Two consecutive physical interaction paths with exact MPS ground
    realizations can be concatenated when their ambient bond dimensions
    agree and their endpoint tensors coincide under that identification.
    The resulting continuous tensor path realizes the ground lines of
    the concatenated physical interaction path.
\end{definition}

\begin{proof}\leanok
    \uses{thm:spt_common_space_path_concatenation}
    Use the first tensor family on the first half of the interval and
    the second on the second half, with the same affine parameter changes
    as the physical path. Endpoint equality gives continuity at the join.
    Injective representatives, nonvanishing, and the ground-line identity
    follow from the corresponding half of the path.
\end{proof}

\begin{lemma}[One-site paths give the independent exact MPS phase condition]
    \label{thm:spt_exact_phase_one_site_path}
    \lean{MPSTensor.blockKron_one_eq_submatrix,
        MPSTensor.rotatePhysical_submatrix_singleBlockEquiv_blockTensor_one,
        MPSTensor.isSameExactMPSGappedPhase_of_oneSite_path}
    \leanok
    \uses{def:spt_exact_mps_gapped_phase}
    Suppose the full endpoint physical representations have orthogonal
    isometric inclusions in a common physical representation, intertwining
    the on-site actions. A symmetric gapped path between their included
    canonical parents, with an exact MPS ground realization whose endpoint
    vector lines are the included endpoint lines, proves the independent
    exact MPS phase condition with blocking length one and trivial
    endpoint characters.
\end{lemma}

\begin{proof}\leanok
    Identify a one-letter word with its single physical index. The blocked
    action is the original action under this identification, and inclusion
    of the blocked tensor gives exactly the original included tensor.
    Isometry, orthogonality, and covariance are preserved by the same
    reindexing. The given path and ground realization therefore provide
    all the required data.
\end{proof}

\begin{lemma}[Unitary preparation of a virtual implementation]
    \label{thm:spt_exact_phase_unitary_preparation}
    \lean{MPSTensor.rotatePhysical_covariance_of_unitary_virtual_action,
        MPSTensor.exists_unitary_covariant_preparation_with_virtual_class}
    \leanok
    Let an injective tensor of positive bond dimension have an
    invertible virtual projective implementation of its exact physical
    symmetry. A nonzero scalar rescaling and a bond change of basis give
    an injective tensor with a unitary virtual implementation in the
    same cohomology class. Its physical covariance is adjoint conjugation.
\end{lemma}

\begin{proof}\leanok
    \uses{thm:spt_prepared_polar_virtual_class}
    The virtual implementation gives invariance of every periodic
    vector. Apply the prepared polar construction preserving its virtual
    class. At its injective endpoint the virtual matrices are unitary;
    their inverses are therefore their adjoints.
\end{proof}

\begin{definition}[Three-stage exact ground realization]
    \label{def:spt_exact_ground_reverse_three_stage}
    \lean{MPSTensor.ExactMPSGroundPath.reverseTransTrans}
    \leanok
    \uses{def:spt_exact_ground_path_concatenation}
    Three exact MPS ground realizations with a common ambient bond
    dimension can be joined after reversing the first path, provided
    their tensors agree at the two joins. This gives the continuous
    ground realization of the same three-stage physical path.
\end{definition}

\begin{lemma}[Dimensions and endpoints of exact ground concatenation]
    \label{thm:spt_exact_ground_three_stage_endpoints}
    \lean{MPSTensor.ExactMPSGroundPath.trans_bondDimension,
        MPSTensor.ExactMPSGroundPath.reverseTransTrans_tensor_zero,
        MPSTensor.ExactMPSGroundPath.reverseTransTrans_tensor_one}
    \leanok
    \uses{def:spt_exact_ground_reverse_three_stage}
    Concatenation retains the common ambient bond dimension. The initial
    tensor of the three-stage realization is the final tensor of the
    first, unreversed family, and its final tensor is the final tensor
    of the third family.
\end{lemma}

\begin{proof}\leanok
    \uses{thm:spt_common_physical_canonical_ground_path}
    Apply the endpoint identities for two successive concatenations;
    reversal exchanges the endpoints of the first family.
\end{proof}

\begin{definition}[Exact polar ground families in the common bond space]
    \label{def:spt_common_physical_polar_ground_paths}
    \lean{MPSTensor.commonPhysicalLeftPolarExactMPSGroundPath,
        MPSTensor.commonPhysicalRightPolarExactMPSGroundPath}
    \leanok
    \uses{def:spt_common_physical_polar_paths, def:spt_supported_bond_lift}
    The two common-space polar parent paths have exact MPS ground
    realizations in the common bond space of dimension $D_0+D_1$.
    For endpoint $j$, include the scaled polar tensor family by the
    corresponding coordinate isometry and take its supported lift
    with initial tensor the included weighted fixed-point endpoint.
\end{definition}

\begin{proof}\leanok
    \uses{thm:spt_canonical_injective_ground_path,
        thm:spt_supported_bond_lift,
        thm:spt_common_physical_polar_endpoints}
    The small-bond polar tensors are continuous and one-site injective,
    and their canonical parents are the prescribed physical interactions.
    The weighted endpoint satisfies the bond intertwining and row-support
    identities. Its supported lift preserves all positive-length
    periodic vectors and therefore the same exact ground realization.
\end{proof}

\begin{lemma}[Exact endpoints of the common polar ground families]
    \label{thm:spt_common_physical_polar_ground_endpoints}
    \lean{MPSTensor.commonPhysicalLeftPolarExactMPSGroundPath_tensor_zero,
        MPSTensor.commonPhysicalRightPolarExactMPSGroundPath_tensor_zero,
        MPSTensor.commonPhysicalLeftPolarExactMPSGroundPath_samePositiveMpvRay_one,
        MPSTensor.commonPhysicalRightPolarExactMPSGroundPath_samePositiveMpvRay_one}
    \leanok
    \uses{def:spt_common_physical_polar_ground_paths}
    At parameter zero, each polar ground family is exactly the included
    weighted fixed-point tensor. At parameter one, its positive-length
    periodic vector lines are those of the corresponding included
    original tensor. Thus both fixed-point joins are exact tensor
    identities, while the original endpoints have the required vector
    lines.
\end{lemma}

\begin{proof}\leanok
    \uses{thm:spt_supported_bond_lift}
    A supported lift has its prescribed initial tensor. At the other
    endpoint its bond intertwiner preserves every nonempty word trace;
    the fixed nonzero normalization changes only the vector line.
\end{proof}

\begin{definition}[Exact MPS realization of the full common-space path]
    \label{def:spt_common_physical_injective_ground_path}
    \lean{MPSTensor.commonPhysicalInjectiveExactMPSGroundPath}
    \leanok
    \uses{def:spt_exact_mps_gapped_phase}
    For two one-site injective tensors of positive bond dimension with
    unitary physical and virtual actions giving exact adjoint covariance
    and the same virtual factor system, the constructed common-space
    path has a continuous exact MPS ground realization. It joins the
    two included endpoint vector lines through the polar deformations
    and the canonical fixed-point comparisons.
\end{definition}

\begin{proof}\leanok
    \uses{thm:spt_supported_bond_lift,
        thm:spt_common_physical_polar_ground_endpoints,
        thm:spt_common_physical_canonical_ground_path,
        thm:spt_common_physical_injective_path,
        def:spt_exact_ground_reverse_three_stage}
    The supported bond lifts preserve the polar periodic vectors and
    give exactly the weighted fixed-point endpoint tensors. They
    therefore join the canonical fixed-point ground realization.
    Reverse the first polar family and concatenate the three families
    with the same parameter changes as the physical path.
\end{proof}

\begin{lemma}[Cohomologous virtual factors give the independent exact MPS phase]
    \label{thm:spt_exact_phase_forward}
    \lean{MPSTensor.isSameExactMPSGappedPhase_of_unitary_common_factor,
        MPSTensor.isSameExactMPSGappedPhase_of_unitary_cohomologous,
        MPSTensor.isSameExactMPSGappedPhase_of_isInjective_cohomologous}
    \leanok
    \uses{def:spt_exact_mps_gapped_phase}
    Let $A_0,A_1$ be one-site injective tensors of positive bond dimension
    with unitary physical actions. Suppose invertible virtual
    projective representations implement their exact tensor symmetries.
    If their factor systems are cohomologous, the tensors satisfy the
    independent exact MPS phase condition. The whole endpoint physical
    spaces have orthogonal isometric images in the constructed common
    physical representation. No virtual family or physical path is
    assumed. This proves the forward construction of
    \cite[Section~II.F.2]{Schuch2011Phases} in the stated one-site,
    trivial-character setting.
\end{lemma}

\begin{proof}\leanok
    \uses{def:spt_common_physical_injective_ground_path,
        thm:spt_common_physical_endpoint_orthogonality,
        thm:spt_exact_phase_one_site_path,
        thm:spt_unitary_factor_alignment,
        thm:spt_exact_phase_unitary_preparation,
        thm:spt_exact_phase_gauge_invariance}
    For unitary virtual actions with a common factor, use the constructed
    common-space path and its exact ground realization. The physical
    inclusions are isometric, orthogonal, and intertwine the original
    endpoint actions. One-site blocking gives the phase condition.
    A unit-circle rephasing aligns cohomologous unitary factors without
    changing their adjoint covariance. For general invertible virtual
    actions, prepare both endpoints with unitary virtual actions in
    the same classes, apply the unitary construction, and transfer
    the phase condition back by gauge invariance.
\end{proof}

\begin{lemma}[Chain characters from an independent physical path]
    \label{thm:spt_physical_path_chain_character}
    \lean{Matrix.eigenvector_of_mem_span_singleton,
        Matrix.eigenvalue_constant_of_continuous_invariant_line,
        Matrix.exists_constant_character_of_continuous_invariant_line,
        MPSTensor.mem_span_mpv_of_commute_ground_line,
        MPSTensor.ExactMPSGroundPath.exists_chain_character,
        MPSTensor.ExactMPSGroundPath.exists_chain_character_mpv}
    \leanok
    \uses{def:spt_exact_mps_gapped_phase}
    For an independent exact MPS ground-state path and every $N\ge2$,
    there is one unit-modulus character $c_N$ such that the physical
    symmetry acts on every length-$N$ ground vector by $c_N(g)$,
    independently of the path parameter. The ambient tensors need not
    be injective, and their minimal bond dimensions may change.
\end{lemma}

\begin{proof}\leanok
    The nonzero periodic vectors depend continuously on the ambient tensors.
    Commutation with the Hamiltonian preserves its unique ground line.
    The eigenvalue on this line depends continuously on the vector and
    belongs to the fixed finite symmetry spectrum; it is therefore
    constant. Multiplication of symmetry operators gives the character
    law, and unitarity gives unit modulus.
\end{proof}

\begin{lemma}[A common physical character with varying minimal dimension]
    \label{thm:spt_varying_bond_physical_character}
    \lean{LinearMap.map_eq_smul_of_span_singleton_eq,
        MPSTensor.mpv_twistedTensor_eq_finKronecker_mulVec,
        MPSTensor.SamePositiveMpvRay.symmetry_eigenvalue,
        MPSTensor.ExactMPSGroundPath.exists_single_character,
        MPSTensor.ExactMPSGroundPath.exists_exact_symmetry}
    \leanok

    An independent exact MPS ground-state path has one unit-modulus
    character $\chi$ such that its physical symmetry acts on every
    positive-length periodic vector by $\chi(g)^N$. Rephasing the physical
    representation by $\chi^{-1}$ makes every ambient tensor exactly
    symmetric, even when its minimal injective bond dimension varies.
    No continuous choice of minimal tensors is assumed.
    This removes the scalar character in the varying-dimension argument
    of \cite[Appendices~B and~C]{Schuch2011Phases}; continuity of the
    canonical virtual support remains to be established.
\end{lemma}

\begin{proof}\leanok
    \uses{thm:spt_physical_path_chain_character,
        thm:spt_character_from_positive_lengths,
        thm:spt_physical_character_twist}
    Equality of nonzero vector rays transfers symmetry eigenvalues to
    each pointwise minimal injective representative. The characters
    $c_2$ and $c_3$ of the fixed finite-chain ground lines are independent
    of the path parameter. The ratio $c_3/c_2$ determines the same one-site
    character for each minimal representative, irrespective of dimension.
    The fundamental theorem gives its length-one symmetry as well.
    Transfer these identities back through equality of periodic rays,
    and remove the common character from the physical action.
\end{proof}

\begin{lemma}[The physical converse with a fixed injective realization]
    \label{thm:spt_fixed_bond_physical_converse}
    \lean{MPSTensor.ExactMPSGroundPath.exists_single_character_of_isInjective,
        MPSTensor.ExactMPSGroundPath.exists_exact_symmetry_of_isInjective,
        MPSTensor.ExactMPSGroundPath.exists_continuous_virtual_data_of_isInjective,
        MPSTensor.ExactMPSGroundPath.exists_rephasing_endpoint_cohomology_of_isInjective}
    \leanok

    Suppose an independent exact MPS ground-state path has a continuous
    ambient tensor realization that is one-site injective at every
    parameter. Then one fixed unit-modulus physical character makes
    the entire family exactly symmetric. Continuous invertible virtual
    projective matrices and continuous factor systems follow, and the
    endpoint factor systems are cohomologous. The endpoint classes agree
    for any representatives of that common rephased symmetry.
    Neither exact tensor symmetry nor a virtual path is an assumption.
    The fixed injective realization is the scope restriction; the
    varying-support converse remains open
    \cite{gap:spc11_spt_interpolation_upper_range}.
\end{lemma}

\begin{proof}\leanok
    \uses{thm:spt_physical_path_chain_character,
        thm:spt_character_from_positive_lengths,
        thm:spt_physical_character_twist,
        thm:spt_tensor_path_class_invariance}
    The chain characters are independent of the path parameter.
    Their ratio $c_3/c_2$ gives one site character for every tensor.
    Remove it from the physical action, lift the resulting exact tensor
    symmetry to continuous virtual matrices, and apply the determinant
    argument. Gauge uniqueness compares arbitrary endpoint choices.
\end{proof}

\begin{definition}[Schmidt coefficient matrix of the two-site ring]
    \label{def:spt_two_site_schmidt_matrix}
    \lean{MPSTensor.twoSiteSchmidtMatrix}
    \leanok
    The coefficient matrix across the one-site cut of the periodic
    two-site MPS is $C_2(A)_{ij}=\operatorname{tr}(A^iA^j)$.
\end{definition}

\begin{lemma}[Minimal bond dimension is locally bounded below]
    \label{thm:spt_minimal_bond_dimension_lower_bound}
    \lean{MPSTensor.rank_twoSiteSchmidtMatrix_of_isInjective,
        MPSTensor.twoSiteSchmidtMatrix_apply_eq_mpv,
        MPSTensor.continuous_twoSiteSchmidtMatrix,
        MPSTensor.SamePositiveMpvRay.rank_twoSiteSchmidtMatrix_eq,
        MPSTensor.ExactMPSGroundPath.eventually_minimalBondDimension_ge}
    \leanok
    \uses{def:spt_two_site_schmidt_matrix, def:spt_exact_mps_gapped_phase}
    A one-site injective tensor of bond dimension $D$ has
    $\operatorname{rank}C_2(A)=D^2$. Along a continuous exact MPS ground
    realization, suppose an injective representative at $t_0$ has
    positive bond dimension $D_0$. In a neighborhood of $t_0$ in the
    parameter interval, every injective representative of the same
    positive-length periodic vector lines has bond dimension at least
    $D_0$. This proves the local dimension-direction assertion in
    \cite[Section~II.F.2, lines~1059--1068]{Schuch2011Phases}.
    It uses the two-site vector and its continuity, and does not derive
    a surviving single-bond projective subspace from the physical gap.
\end{lemma}

\begin{proof}\leanok
    \uses{thm:ldpolar_injective_physical}
    With $M$ the physical matrix and $W$ the swap of its two virtual
    indices, $C_2=MWM^T$. Injectivity gives a left inverse of $M$,
    so multiplication by that inverse and its transpose recovers $W$.
    The two rank inequalities therefore give rank $D^2$.
    Nonzero ray equality rescales $C_2$ by a nonzero scalar and preserves
    its rank. The coefficient matrix depends continuously on the
    ground tensor, and matrix rank is lower semicontinuous. Hence nearby
    representative dimensions satisfy $D^2\ge D_0^2$, or $D\ge D_0$.
\end{proof}

\begin{lemma}[Two-site trace data determine the physical quotient]
    \label{thm:spt_finite_two_site_trace_quotient}
    \lean{MPSTensor.traceGram_eq_smul_of_twoSite_trace_eq,
        MPSTensor.ker_linearCombination_eq_of_twoSite_trace_eq,
        MPSTensor.bondDimension_eq_of_twoSite_trace_eq}
    \leanok
    Let $A^i\in M_D(\C)$ and $B^i\in M_E(\C)$ be one-site injective.
    Suppose $c_2\ne0$ and
    $\tr(B^iB^j)=c_2\tr(A^iA^j)$ for all physical indices $i,j$.
    Then the coefficient maps $x\mapsto\sum_i x_iA^i$ and
    $x\mapsto\sum_i x_iB^i$ have the same kernel, and $D=E$.
    Their trace forms are the compositions of these coefficient maps
    with the corresponding nondegenerate trace pairings.
\end{lemma}

\begin{proof}\leanok
    Injectivity makes $M\mapsto(\tr(MA^i))_i$ and its analogue for $B$
    injective. The proportionality of the two-site trace forms therefore
    identifies the coefficient kernels. Both coefficient maps are
    surjective, so rank--nullity gives $D^2=E^2$ and hence $D=E$.
    The trace forms here are complex bilinear, rather than positive
    Hermitian forms.
\end{proof}

\begin{lemma}[Three-site trace data determine multiplication]
    \label{thm:spt_finite_three_site_trace_algebra}
    \lean{MPSTensor.linearExtension_exists_unique_of_twoSite_trace_eq,
        MPSTensor.linearExtension_mul_of_two_three_trace_eq}
    \leanok
    Let $A,B$ be one-site injective and suppose their two-site traces
    satisfy the preceding proportionality, allowing $c_2=0$ for the
    purely linear assertion. There is a unique linear map
    $T:M_D(\C)\to M_E(\C)$ with $T(A^i)=B^i$.
    If $c_3\ne0$ and
    $\tr(B^iB^jB^k)=c_3\tr(A^iA^jA^k)$ for all $i,j,k$, then
    $T(MN)=(c_2/c_3)T(M)T(N)$.
\end{lemma}

\begin{proof}\leanok
    Use a left inverse of the trace-pairing map for $B$ to extend the
    letter assignment. The letters of $A$ span the full matrix algebra,
    giving uniqueness. Three-site traces prove the multiplication identity
    on pairs of letters. Bilinearity extends it to arbitrary matrices.
\end{proof}

\begin{lemma}[Reconstruction from two finite rings]
    \label{thm:spt_finite_ring_reconstruction}
    \lean{MPSTensor.gaugeEquiv_smul_of_two_three_trace_eq,
        MPSTensor.samePositiveMpvRay_of_two_three_trace_eq,
        MPSTensor.samePositiveMpvRay_of_mpv_two_three_eq_smul}
    \leanok
    \uses{def:gauge_equiv}
    Suppose two one-site injective tensors have proportional periodic
    vectors at lengths two and three, with nonzero constants $c_2,c_3$.
    They have the same minimal bond dimension, and $B$ is related by
    a bond conjugation to $(c_3/c_2)A$. Consequently their periodic
    vector rays agree at every positive length.
    This finite-data assertion does not construct continuous canonical
    coordinates through a changing bond dimension.
\end{lemma}

\begin{proof}\leanok
    \uses{thm:spt_finite_two_site_trace_quotient,
        thm:spt_finite_three_site_trace_algebra}
    The dimension equality identifies the matrix spaces.
    At positive bond dimension, the normalized map $(c_2/c_3)T$ is
    a nonzero multiplicative linear endomorphism. Skolem--Noether represents it by conjugation with an
    invertible matrix. Traces of words are invariant under conjugation;
    multiplying every letter by $c_3/c_2$ rescales the length-$N$ vector
    by its nonzero $N$th power. The zero-dimensional case is immediate.
\end{proof}

\begin{definition}[Coordinates for the trace-quotient product]
    \label{def:spt_trace_quotient_product_coordinates}
    \lean{Matrix.traceQuotientProductCoordinates}
    \leanok
    Let $G\in M_d(\C)$ be the complex-bilinear two-site trace form,
    let $F:\C^r\to\C^d$ be a coefficient section, and put $C=GF$.
    For contracted three-site columns $R:\C^{r\times r}\to\C^d$, set
    \begin{align}
        M        & =(C^\dagger C)^{-1}C^\dagger R,
        \notag                                     \\
        \mu(x,y) & =M(x_a y_b)_{(a,b)}.
        \notag
    \end{align}
    The positive Hermitian matrix $C^\dagger C$ serves to invert the
    section. No positivity of the bilinear trace form $G$ is assumed.
\end{definition}

\begin{lemma}[A fixed section of a constant-rank trace form]
    \label{thm:spt_trace_quotient_fixed_section}
    \lean{Matrix.exists_fixedSection_of_rank,
        Matrix.exists_open_fixedSection_range}
    \leanok
    Let $G(t)$ be a continuous matrix family of constant rank $r$.
    At any base point $t_0$ there is a coefficient section $F$ such
    that $G(t_0)F$ is injective. For this fixed $F$, there is an open
    neighborhood $S$ of $t_0$ on which $G(t)F$ remains injective and
    has the same range as $G(t)$.
\end{lemma}

\begin{proof}\leanok
    Choose coordinates on the range of $G(t_0)$ and lift them through
    $G(t_0)$. The Hermitian Gram matrix of the resulting section is
    invertible. Invertibility persists on an open neighborhood.
    The section range is contained in the full range and both have
    dimension $r$, so they coincide there.
\end{proof}



\begin{lemma}[Continuous multiplication coefficients at constant rank]
    \label{thm:spt_trace_quotient_continuous_product}
    \lean{Matrix.continuous_traceQuotientProductCoordinates,
        Matrix.exists_open_continuous_traceQuotientProductCoordinates,
        Matrix.exists_local_continuous_traceQuotientProductCoordinates}
    \leanok
    \uses{def:spt_trace_quotient_product_coordinates,
        thm:spt_trace_quotient_fixed_section}
    Suppose $G(t)$ has constant rank $r$, and for each fixed section
    $F$ the contracted columns $R_F(t)$ depend continuously on $t$.
    Then $F$ can be chosen so that the multiplication coefficients
    $(C(t)^\dagger C(t))^{-1}C(t)^\dagger R_F(t)$ are continuous on
    an open neighborhood of the base point. The conclusion includes
    the section and its open neighborhood.
\end{lemma}

\begin{proof}\leanok
    Use the preceding fixed section. Its Gram matrix is positive
    definite throughout the neighborhood. Matrix inversion and
    multiplication are continuous there.
\end{proof}

\begin{definition}[Contracted three-site trace columns]
    \label{def:spt_trace_quotient_triple_columns}
    \lean{Matrix.traceQuotientTripleColumns}
    \leanok
    For an ambient tensor $B$ and a coefficient section $F$, define
    \begin{align}
        R_F(B)_{j,(a,b)}
         & =\tr\left[
            \left(\sum_i F_{ia}B^i\right)
            \left(\sum_k F_{kb}B^k\right)B^j
            \right].
        \notag
    \end{align}
\end{definition}

\begin{lemma}[Contraction of finite-ring trace data]
    \label{thm:spt_trace_quotient_contraction}
    \lean{Matrix.traceQuotientSection_pair,
        Matrix.traceQuotientSection_triple,
        Matrix.traceQuotientProductColumns_triple,
        Matrix.continuous_traceQuotientTripleColumns}
    \leanok
    \uses{def:spt_trace_quotient_triple_columns}
    Write $Qx=\sum_i(Fx)_i A^i$ and
    $\Phi_A(X)=(\tr(XA^j))_j$.
    If the two- and three-site traces of $B$ are proportional to those
    of $A$, with constants $\alpha_2,\alpha_3$, then
    \begin{align}
        GFx                     & =\alpha_2\Phi_A(Qx),
        \notag                                             \\
        R_F(B)(x_a y_b)_{(a,b)} & =\alpha_3\Phi_A(Qx\,Qy).
        \notag
    \end{align}
    The contracted columns depend continuously on a continuous ambient
    tensor family in a fixed bond dimension.
\end{lemma}

\begin{proof}\leanok
    Expand the coefficient sums and use bilinearity and trilinearity
    of the trace expressions. Continuity follows from the finite sums
    and matrix products in their defining formula.
\end{proof}

\begin{lemma}[Recovery of the minimal matrix algebra]
    \label{thm:spt_trace_quotient_matrix_algebra_recovery}
    \lean{Matrix.traceQuotientProductCoordinates_recover,
        Matrix.exists_linearEquiv_traceQuotientSection,
        Matrix.traceQuotientProductCoordinates_algebraRecovery,
        Matrix.traceQuotientProductCoordinates_algebraRecovery_of_two_three_trace_eq}
    \leanok
    \uses{def:spt_trace_quotient_product_coordinates,
        thm:spt_trace_quotient_contraction}
    Suppose $r=D^2$, $GF$ is injective, and the preceding trace
    proportionalities hold with $\alpha_2\alpha_3\ne0$.
    Then $Q:\C^r\to M_D(\C)$ is a linear isomorphism and
    \begin{align}
        Q\mu(x,y) & =\frac{\alpha_3}{\alpha_2}Qx\,Qy.
        \notag
    \end{align}
    Consequently $E=(\alpha_3/\alpha_2)Q$ identifies $\mu$ with
    matrix multiplication, and $\mu$ is associative.
    The minimal tensor and the identification are used only pointwise;
    neither is assumed continuous. Together with the preceding
    continuity result, this gives continuous multiplication coefficients
    at constant trace rank. Choosing continuous tensor matrices,
    particularly through a change of rank, remains a separate step in
    the reconstruction proposed in
    \cite[Section~II.F.2, lines~953--993]{Schuch2011Phases}.
\end{lemma}

\begin{proof}\leanok
    Injectivity of $GF$ implies injectivity of $Q$; the dimension
    equality gives surjectivity. The left inverse
    $(C^\dagger C)^{-1}C^\dagger$ cancels $C$ in the contracted
    three-site identity. Rescaling $Q$ by $\alpha_3/\alpha_2$ makes
    the resulting identification multiplicative. Associativity follows
    from associativity of matrix multiplication and injectivity of $E$.
\end{proof}

\begin{lemma}[Local continuous algebra from an ambient tensor family]
    \label{thm:spt_trace_quotient_continuous_algebra}
    \lean{Matrix.exists_local_continuous_traceQuotientAlgebra}
    \leanok
    \uses{thm:spt_trace_quotient_continuous_product,
        thm:spt_trace_quotient_matrix_algebra_recovery}
    Let $B(t)$ be continuous in a fixed ambient bond dimension.
    Suppose at each parameter it has a one-site injective representative
    $A(t)$ whose two- and three-site traces are proportional to those
    of $B(t)$ by nonzero constants $\alpha_2(t),\alpha_3(t)$.
    Assume $D(t)^2=r$ is independent of $t$.
    No continuity of $A(t)$ or the proportionality constants is required.
    Near any base point, there is a fixed coefficient section for which
    the multiplication coefficients are continuous and the product is
    associative. At each parameter its normalized identification with
    $M_{D(t)}(\C)$ is
    \begin{align}
        E_t x & =\frac{\alpha_3(t)}{\alpha_2(t)}
        \sum_i(Fx)_iA^i(t).
        \notag
    \end{align}
    This constant-dimension assertion does not construct a continuous
    tensor realization through a change of rank
    \cite{gap:spc11_spt_interpolation_upper_range}.
\end{lemma}

\begin{proof}\leanok
    The coefficient map of $A(t)$ is surjective, and its trace-pairing
    map is injective. Their composition, multiplied by the nonzero
    scalar $\alpha_2(t)$, therefore has rank $D(t)^2=r$.
    Apply the continuous-section construction and the pointwise
    matrix-algebra recovery on the resulting neighborhood.
\end{proof}

\begin{lemma}[A continuous unit for the local trace quotient]
    \label{thm:spt_trace_quotient_continuous_unit}
    \lean{Matrix.exists_local_continuous_traceQuotientUnit}
    \leanok
    \uses{thm:spt_trace_quotient_continuous_algebra}
    Under the hypotheses of the preceding theorem, the section and
    neighborhood can be chosen so that the quotient product also has
    a continuous two-sided unit $u(t)$. The same pointwise normalized
    identifications with the minimal matrix algebra remain available.
    No continuous unit is assumed.
\end{lemma}

\begin{proof}\leanok
    At the base point choose the unit $u_0$ using its matrix-algebra
    identification. Left multiplication $L_t(u_0)$ is the identity at
    the base point, and is invertible on a smaller open neighborhood.
    Define $u(t)=L_t(u_0)^{-1}u_0$. Any pointwise two-sided unit satisfies
    $L_t(u_0)u(t)=u_0$ and hence equals this continuous choice.
\end{proof}

\begin{definition}[Repeated bond copies of a product tensor]
    \label{def:spt_repeated_product_tensor}
    \lean{MPSTensor.repeatedProductTensor}
    \leanok
    Set $R_D^0=I_D$ and $R_D^1=0$.
\end{definition}

\begin{lemma}[Physical rays do not fix an ambient stationary space]
    \label{thm:spt_redundant_bond_stationary_space}
    \lean{MPSTensor.repeatedProductTensor_adjointMapLM,
        MPSTensor.repeatedProductTensor_isUnital,
        MPSTensor.repeatedProductTensor_fixedSpace_finrank,
        MPSTensor.repeatedProductTensor_mpv,
        MPSTensor.repeatedProductTensor_samePositiveMpvRay}
    \leanok
    \uses{def:spt_repeated_product_tensor}
    The tensor $R_D$ is unital and its adjoint transfer map is the
    identity on $M_D(\C)$. Its stationary space therefore has complex
    dimension four when $D=2$. Yet the periodic vectors of $R_2$ are
    twice those of $R_1$, so their positive-length rays agree.
    In particular, the physical rays alone do not determine the
    stationary-space dimension of an arbitrary ambient realization.
\end{lemma}

\begin{proof}\leanok
    Only the zero physical letter contributes to either the transfer map
    or a nonvanishing periodic word. The transfer map fixes every matrix,
    while a nonvanishing periodic coefficient is the trace of $I_D$, namely
    $D$. Thus the two realizations differ by the constant factor two.
\end{proof}

\begin{lemma}[Local continuous coordinates for a fixed-rank support]
    \label{thm:spt_continuous_projection_frame}
    \lean{Matrix.continuous_polarPos_family,
        Matrix.polarIso_eq_mul_inv_polarPos_of_injective,
        Matrix.continuous_polarIso_family_of_injective,
        Matrix.exists_open_projected_frame_injective,
        Matrix.exists_local_continuous_isometric_projected_frame,
        Matrix.range_eq_of_isometric_frame_of_rank_eq,
        Matrix.frame_mul_conjTranspose_eq_of_rank_eq,
        Matrix.exists_local_continuous_projection_frame}
    \leanok
    Let $P(t)$ be a continuous family of orthogonal projections of rank $r$,
    and let $J_0$ be an isometric frame with $J_0J_0^\dagger=P(t_0)$.
    Near $t_0$, the polar isometry $J(t)$ of $P(t)J_0$ depends continuously
    on $t$ and satisfies $J(t)^\dagger J(t)=I$ and
    $J(t)J(t)^\dagger=P(t)$. The positive polar factor is continuous for
    every continuous rectangular matrix family; its polar isometry is
    continuous wherever the matrices are injective.
    These are auxiliary ingredients for the support argument in
    \cite[Appendix~C]{Schuch2011Phases}. Continuity and fixed rank of $P$
    are hypotheses here; deriving the relevant canonical supports from
    the physical gap remains open
    \cite{gap:spc11_spt_interpolation_upper_range}.
\end{lemma}

\begin{proof}\leanok
    The Gram matrix of $P(t)J_0$ equals the identity at $t_0$ and remains
    invertible nearby. Continuous functional calculus gives continuity
    of its positive square root, and inversion is continuous on the
    invertible matrices. Thus the polar isometry is continuous.
    Its range lies in that of $P(t)$ and has the same dimension, so the
    two orthogonal range projections coincide.
\end{proof}

\begin{lemma}[Continuous separated spectral corners]
    \label{thm:spt_continuous_spectral_corner}
    \lean{Matrix.spectralCornerRamp,
        Matrix.continuous_spectralCornerRamp,
        Matrix.spectralCornerRamp_eq_zero_of_le,
        Matrix.spectralCornerRamp_eq_one_of_le,
        Matrix.spectralCorner,
        Matrix.isHermitian_spectralCorner,
        Matrix.spectralCorner_mul_self,
        Matrix.continuous_spectralCorner_family,
        Matrix.spectralCorner_eq_cfc_highIndicator,
        Matrix.commute_spectralCorner_of_commute,
        Matrix.unitary_conj_spectralCorner_eq,
        Matrix.isLocallyConstant_rank_of_continuous_projection_family,
        Matrix.isLocallyConstant_rank_spectralCorner_family}
    \leanok
    Let $A(t)$ be a continuous Hermitian matrix family and let $a<b$.
    Suppose every spectrum is contained in
    $(-\infty,a]\cup[b,\infty)$. A continuous scalar function equal to
    zero below $a$ and one above $b$ gives, by functional calculus, the
    orthogonal spectral projection $P(t)$ onto the part above the cutoff.
    This projection is continuous and has locally constant rank.
    Every operator commuting with $A(t)$ commutes with $P(t)$; conjugation
    by a commuting unitary fixes the corner projection.
    The full density rank may change below the cutoff.
    These facts provide auxiliary separated supports for the argument in
    \cite[Appendix~C]{Schuch2011Phases}; continuity and spectral separation
    of the relevant density remain explicit hypotheses.
\end{lemma}

\begin{proof}\leanok
    The continuous ramp equals the spectral indicator on every spectrum,
    so its functional calculus is Hermitian and idempotent.
    Continuous functional calculus gives continuity, commutation, and
    invariance under a commuting unitary. The trace of an idempotent
    equals its integer rank. Its continuous trace therefore makes the
    rank locally constant.
\end{proof}

\begin{lemma}[Continuous invariant compression preserves the factor]
    \label{thm:spt_continuous_invariant_compression}
    \lean{TNLean.Algebra.ProjectiveRepresentation.exists_local_continuous_unitary_compression,
        TNLean.Algebra.ProjectiveRepresentation.continuous_factor_family}
    \leanok

    Suppose unitary projective matrices, continuous in the parameter for
    each group element, act
    on one fixed ambient space, and commutes with a continuous fixed-rank
    family of orthogonal projections of nonzero rank. Near a prescribed
    parameter, continuous isometric frames give continuous unitary
    compressed projective matrices with exactly the original factor system.
    A family of invertible projective matrices, continuous in the parameter
    for each group element, on a positive
    fixed dimension also has continuous scalar factors.
    The continuous ambient representation and projections are explicit
    auxiliary hypotheses for the argument in
    \cite[Appendix~C]{Schuch2011Phases}.
\end{lemma}

\begin{proof}\leanok
    \uses{thm:spt_continuous_projection_frame,
        thm:spt_invariant_bond_compression}
    With $J(t)J(t)^\dagger=P(t)$, compress by
    $J(t)^\dagger V_g(t)J(t)$. Invariance inserts the projection between
    products, preserving the projective multiplication law and unitarity.
    Matrix multiplication gives continuity. To recover the scalar factor,
    take one diagonal entry of
    $V_gV_h\operatorname{adj}(V_{gh})$ and divide by
    $\det(V_{gh})$, which never vanishes.
\end{proof}

\begin{lemma}[Cohomology from a continuous injective tensor path]
    \label{thm:spt_tensor_path_class_invariance}
    \lean{MPSTensor.exists_continuous_virtualGauge_of_isOnSiteSymmetric,
        MPSTensor.exists_continuous_projectiveRepresentation_of_isOnSiteSymmetric,
        MPSTensor.cohomologousTo_of_continuous_isOnSiteSymmetric_tensorPath}
    \leanok

    Let $A_t$ be a continuous family of one-site injective tensors
    of fixed positive bond dimension, with exact on-site symmetry
    under the same action of a group $G$.
    Its virtual symmetry matrices can be chosen invertibly and
    continuously in $t$, separately for each $g\in G$.
    They give projective representations whose factor systems depend
    continuously on $t$ for each pair $g,h$.
    Any virtual representatives of the endpoint tensors have
    cohomologous factor systems. No continuous virtual family is an
    input, and no finiteness assumption on $G$ is needed.
    This is the fixed-dimension, exact-symmetry part of
    \cite[Section~II.F.2]{Schuch2011Phases}; varying minimal support and
    extraction of continuous tensors from an independent ground-state
    path remain open, as recorded in
    \cite{gap:spc11_spt_interpolation_upper_range}.
\end{lemma}

\begin{proof}\leanok
    \uses{thm:spt_global_virtual_gauge,
        thm:spt_virtual_path_class_invariance}
    Apply interval lifting separately to each symmetry twist.
    Uniqueness of the injective bond gauge up to scalar gives the
    projective multiplication law, using inverse-group indexing for
    the virtual action. An entry of the product with an adjugate
    gives continuous factor systems. The determinant argument proves
    their endpoint cohomology, and gauge uniqueness identifies that
    class with any prescribed endpoint representatives.
\end{proof}

\begin{lemma}[Cohomology along a fixed-dimension virtual path]
    \label{thm:spt_virtual_path_class_invariance}
    \lean{MPSTensor.continuous_rootOfUnity_path_const,
        MPSTensor.exists_continuous_pow_root_path,
        MPSTensor.continuous_projectiveFactor_of_unitaryMatrixPath,
        MPSTensor.continuous_projectiveFactor_of_invertibleMatrixPath,
        MPSTensor.projectiveFactor_endpoints_cohomologous_of_det_path,
        MPSTensor.unitary_projectivePath_factor_endpoints_cohomologous,
        MPSTensor.invertible_projectivePath_factor_endpoints_cohomologous}
    \leanok
    \uses{def:cohomologous_to}
    Let $G$ be any group and $D>0$. Suppose matrices
    $V_\gamma(g)\in\mathrm{GL}_D(\C)$ depend continuously on
    $\gamma\in[0,1]$ for each $g\in G$ and obey
    $V_\gamma(g)V_\gamma(h)=\omega_\gamma(g,h)V_\gamma(gh)$.
    Then $\omega_0$ and $\omega_1$ are cohomologous as
    $\C^\times$-valued factor systems. In particular, this holds for
    unitary virtual matrices. This theorem assumes the
    virtual path; extracting compatible continuous tensors from an
    independent exact MPS ground-state path, including changes of
    bond dimension, is a separate step in the classification theorem
    \cite[Section~II.F.2, lines~930--946]{Schuch2011Phases}.
\end{lemma}

\begin{proof}\leanok
    Each $\det V_\gamma(g)$ is a continuous nonzero complex path, so it
    has a continuous $D$th root $\beta_\gamma(g)$ on the interval.
    Taking determinants in $V_\gamma(g)V_\gamma(h)=\omega_\gamma(g,h)V_\gamma(gh)$
    gives $\beta_\gamma(g)^D\beta_\gamma(h)^D=\omega_\gamma(g,h)^D\beta_\gamma(gh)^D$,
    so the rescaled factor system
    \begin{align}
        \eta_\gamma(g,h)
         & =\omega_\gamma(g,h)\,\frac{\beta_\gamma(gh)}{\beta_\gamma(g)\beta_\gamma(h)}
         &
        \eta_\gamma(g,h)^D
         & =1
        \notag
    \end{align}
    takes values among the $D$th roots of unity.  It depends continuously on
    $\gamma$, hence is constant, and $\eta_0=\eta_1$ reads
    \begin{align}
        \omega_1(g,h)
         & =\psi(g)\psi(h)\psi(gh)^{-1}\omega_0(g,h),
         &
        \psi(g)
         & =\beta_1(g)/\beta_0(g).
        \notag
    \end{align}
\end{proof}

\begin{lemma}[Restriction to an invariant bond subspace]
    \label{thm:spt_invariant_bond_compression}
    \lean{Matrix.isometry_compression_mul_of_commute,
        Matrix.isometry_compression_mem_unitaryGroup,
        TNLean.Algebra.ProjectiveRepresentation.exists_unitary_compression}
    \leanok
    \uses{def:projective_rep}
    Let $V_g$ be a unitary projective representation on $\C^D$ with
    factor system $\omega$, and let $J:\C^r\to\C^D$ be an isometry,
    where $r>0$. Suppose $V_g$ commutes with $JJ^\dagger$ for every $g$.
    Then $J^\dagger V_gJ$ is a unitary projective representation with
    exactly the same factor system $\omega$.
    This is the algebraic restriction in
    \cite[Appendix~C]{Schuch2011Phases}. The invariant bond subspace
    is given here; deriving it continuously from the physical gap
    remains a separate problem, recorded in
    \cite{gap:spc11_spt_interpolation_upper_range}.
\end{lemma}

\begin{proof}\leanok
    Invariance and $J^\dagger J=I$ give
    $(J^\dagger XJ)(J^\dagger YJ)=J^\dagger XYJ$ whenever
    $X$ preserves the included subspace. Apply this identity to $V_g,V_h$
    for the projective law, and to $V_g,V_g^\dagger$ for unitarity.
\end{proof}

\begin{lemma}[Spectral corners retaining the base support]
    \label{thm:spt_local_spectral_support}
    \lean{Matrix.eventually_spectral_separation_of_continuousAt,
        Matrix.exists_positive_spectral_cutoffs,
        Matrix.spectralCorner_eq_supportProj_of_nonzero_spectrum_separation,
        Matrix.exists_open_constantRank_spectralCorner,
        Matrix.exists_local_continuous_spectral_frame,
        Matrix.supportProj_mul_spectralCorner}
    \leanok
    Let $\sigma(t)$ be a continuous Hermitian matrix family, with
    $\sigma(t_0)\geq 0$, and let $J_0$ be an isometric frame of its support.
    There are $0<a<b$ and a neighborhood of $t_0$ on which the spectral
    projection $P(t)$ onto the eigenvalues above the cutoffs has constant
    rank and satisfies $P(t_0)=\operatorname{supp}\sigma(t_0)$.
    The polar isometry $J(t)$ of $P(t)J_0$ is continuous there, with
    $J(t_0)=J_0$, $J(t)^\dagger J(t)=I$, and
    $J(t)J(t)^\dagger=P(t)$.
    When $\sigma(t)\geq0$, the cutoff always satisfies
    $\operatorname{supp}\sigma(t)\,P(t)=P(t)$, whether or not its spectrum
    is separated at that parameter.
\end{lemma}

\begin{proof}\leanok
    \uses{thm:spt_continuous_spectral_corner,
        thm:spt_continuous_projection_frame}
    The real spectrum at $t_0$ is finite, so its nonzero positive values
    have a positive lower bound. Choose both cutoffs below that bound.
    Upper hemicontinuity of the spectrum preserves the intervening gap
    nearby. Functional calculus gives a continuous orthogonal projection,
    whose integer rank is locally constant. Apply the continuous-frame
    construction to this projection. At $t_0$, $P(t_0)J_0=J_0$, whose
    polar isometry is $J_0$ itself. The cutoff vanishes at zero, so its
    range lies in the support of every positive semidefinite matrix.
\end{proof}

\begin{definition}[Individual matrix extension along an isometry]
    \label{def:spt_supported_matrix_extension}
    \lean{Matrix.isometryExtendByIdentity}
    \leanok
    For an isometry $K:\C^D\to\C^k$ and a matrix $X\in\MN{D}$, set
    $\widetilde X=KXK^\dagger+(I-KK^\dagger)$.
\end{definition}

\begin{lemma}[Comparison of supported bond frames]
    \label{thm:spt_supported_frame_comparison}
    \lean{Matrix.isometryExtendByIdentity_mul_isometry,
        Matrix.isometry_conjTranspose_mul_isometryExtendByIdentity,
        Matrix.isometry_reconstruct_of_supportProj,
        Matrix.commute_isometryExtendByIdentity_of_commute_compression,
        Matrix.isometry_absorption_of_projector_absorption,
        Matrix.isIsometry_conjTranspose_mul_of_absorption,
        Matrix.isometry_double_compression_eq_of_absorption,
        Matrix.commute_compression_of_commute_isometryExtendByIdentity,
        Matrix.supportFrame_dimension_pos_of_trace_one}
    \leanok
    \uses{def:spt_supported_matrix_extension}
    Let $K:\C^D\to\C^k$ be isometric. Then
    $\widetilde X K=KX$ and $K^\dagger\widetilde X=XK^\dagger$.
    If $\sigma\geq0$ has support $KK^\dagger$, then
    $\sigma=K(K^\dagger\sigma K)K^\dagger$; commutation of $X$ with
    $K^\dagger\sigma K$ implies commutation of $\widetilde X$ with $\sigma$.
    If $J:\C^r\to\C^k$ is isometric and $KK^\dagger J=J$, then
    $K^\dagger J$ is isometric and
    $(K^\dagger J)^\dagger(K^\dagger M K)(K^\dagger J)=J^\dagger M J$.
    The absorption $KK^\dagger JJ^\dagger=JJ^\dagger$ already implies
    $KK^\dagger J=J$.
    An operator $P$ commuting with $\widetilde X$ has compression
    $K^\dagger P K$ commuting with $X$.
    Finally, $\tr\sigma=1$ forces $D>0$.
\end{lemma}

\begin{proof}\leanok
    Multiply the defining matrix extension by $K$ and $K^\dagger$ and
    use $K^\dagger K=I$. The support projection absorbs $\sigma$ on
    both sides, giving the reconstruction identity. The same identities
    prove the commutation and nested-compression assertions.
    Multiplication of $KK^\dagger JJ^\dagger=JJ^\dagger$ by $J$ gives
    frame absorption. If $D=0$, the support is zero and hence
    $\sigma=0$, contradicting its trace.
\end{proof}

\begin{lemma}[Local cohomology for a continuous supported density]
    \label{thm:spt_supported_density_class_invariance}
    \lean{MPSTensor.eventually_cohomologousTo_of_continuous_supported_density}
    \leanok
    \uses{def:projective_rep}
    For $t\in[0,1]$, let $B(t)$ be a continuous tensor family of fixed
    bond dimension $k$, and let $\sigma(t)$ be a continuous positive semidefinite matrix
    family of trace one. Choose pointwise isometric support frames
    $K(t):\C^{D(t)}\to\C^k$ and put
    $A^i(t)=K(t)^\dagger B^i(t)K(t)$.
    Suppose the unitary virtual actions of $A(t)$ implement a
    fixed exact physical symmetry and commute with
    $K(t)^\dagger\sigma(t)K(t)$.
    If $A(t_0)$ is one-site injective, the virtual factor systems at
    $t$ and $t_0$ are cohomologous for all sufficiently nearby $t$.
    The dimensions $D(t)$, support frames and virtual actions need no
    continuity, and no projective action is required on the complementary
    bond space.
\end{lemma}

\begin{proof}\leanok
    \uses{thm:spt_local_spectral_support,
        thm:spt_supported_frame_comparison,
        thm:spt_local_invariant_compression}
    Retain the separated spectral corner whose base value is the full
    support of $\sigma(t_0)$, and extend $K(t_0)$ to its continuous
    isometric frame $J(t)$. This corner lies in the support of
    $\sigma(t)$ and is invariant under each individually extended virtual
    matrix. Thus $K(t)^\dagger J(t)$ is an invariant isometric frame in
    the minimal bond space. Its compression is the continuous
    tensor $J(t)^\dagger B(t)J(t)$, which equals $A(t_0)$ at the base.
    Apply local cohomology invariance under pointwise compression.
\end{proof}

\begin{lemma}[Cohomology under pointwise invariant compression]
    \label{thm:spt_pointwise_invariant_compression}
    \lean{Matrix.isometry_compression_conj_of_commute,
        MPSTensor.cohomologousTo_of_continuous_injective_pointwise_compression}
    \leanok
    \uses{def:projective_rep}
    Let $A(t)$ have bond dimension $D(t)$, and let $V_g(t)$ be a unitary
    projective representation with factor $\omega_t$ implementing exact
    covariance for a fixed physical action. Let
    $J(t):\C^r\to\C^{D(t)}$ be isometric, where $r>0$, and suppose
    $V_g(t)$ commutes with $J(t)J(t)^\dagger$.
    If $C^i(t)=J(t)^\dagger A^i(t)J(t)$ is continuous and one-site
    injective for $0\leq t\leq1$, then $\omega_1$ and $\omega_0$ are
    cohomologous. No continuity of $D(t)$, $A(t)$, $J(t)$ or $V_g(t)$
    is required, and the tensor letters need not preserve the included
    bond subspace.
\end{lemma}

\begin{proof}\leanok
    \uses{thm:spt_invariant_bond_compression,
        thm:spt_tensor_path_class_invariance}
    Commutation with $JJ^\dagger$ gives
    $J^\dagger X=(J^\dagger XJ)J^\dagger$ and
    $X^\dagger J=J(J^\dagger XJ)^\dagger$. Hence conjugation descends
    to the compressed tensors. The compressed virtual representation
    has exactly the factor $\omega_t$. Apply cohomology invariance for
    continuous injective tensor families to $C(t)$.
\end{proof}

\begin{lemma}[Local cohomology under invariant compression]
    \label{thm:spt_local_invariant_compression}
    \lean{MPSTensor.eventually_isInjective_of_continuousAt,
        Set.Icc.eventually_forall_convexComb,
        MPSTensor.eventually_cohomologousTo_of_continuousOn_pointwise_compression,
        MPSTensor.eventually_cohomologousTo_of_continuous_pointwise_compression,
        MPSTensor.eventually_cohomologousTo_of_continuousOn_invariant_compression}
    \leanok
    Under the pointwise compression assumptions, it suffices that $C(t)$
    be continuous on a neighborhood of $t_0$ and injective at $t_0$.
    Then $\omega_t$ and $\omega_{t_0}$ are cohomologous for all
    sufficiently nearby $t$. For a fixed ambient bond space,
    continuity of $A(t)$ and $J(t)$ on that neighborhood implies the
    required continuity of $C(t)=J(t)^\dagger A(t)J(t)$.
\end{lemma}

\begin{proof}\leanok
    \uses{thm:spt_pointwise_invariant_compression}
    The Gram matrix of the one-site physical coefficient matrix is
    invertible at $t_0$ and remains invertible nearby. Thus injectivity
    persists locally. Every sufficiently short affine subinterval
    issuing from $t_0$ remains in this neighborhood. Reparametrize it
    by $[0,1]$ and apply the pointwise compression theorem.
\end{proof}

\begin{lemma}[Symmetry of a unique stationary matrix]
    \label{thm:spt_unique_stationary_symmetry}
    \lean{MPSTensor.adjointMap_rotatePhysical_eq_of_unitary,
        MPSTensor.adjointMap_conj_of_unitary_covariance,
        MPSTensor.commute_stationaryMatrix_of_unitary_covariance,
        MPSTensor.virtualAction_commute_stationaryMatrix_of_unique}
    \leanok
    Suppose a unitary physical action satisfies
    $\sum_j u_{ij}A^j=XA^iX^\dagger$, where $X$ is unitary.
    If the adjoint transfer map has a unique trace-one stationary
    matrix $\sigma$, then $X\sigma=\sigma X$.
    This conclusion requires neither invertibility nor positivity of
    $\sigma$. In particular, it applies to a singular stationary density.
\end{lemma}

\begin{proof}\leanok
    Unitary physical rotation leaves the adjoint transfer map unchanged.
    Covariance consequently gives
    $\mathcal E^*(XZX^\dagger)=X\mathcal E^*(Z)X^\dagger$.
    Thus $X\sigma X^\dagger$ is another trace-one stationary matrix.
    Uniqueness gives $X\sigma X^\dagger=\sigma$, which is equivalent
    to commutation with $X$.
\end{proof}

\begin{lemma}[Stationary density under support compression]
    \label{thm:spt_stationary_density_compression}
    \lean{MPSTensor.stationaryMatrix_compression_unique}
    \leanok
    Suppose $\sigma\geq0$ is the unique trace-one stationary matrix
    of the adjoint transfer map of $B$. For an isometric support frame
    $K$, put $A^i=K^\dagger B^iK$. Then $K^\dagger\sigma K$ is the
    unique trace-one stationary matrix of the adjoint transfer map of $A$.
    No unitality assumption is needed for this implication.
\end{lemma}

\begin{proof}\leanok
    Stationarity makes the support of $\sigma$ invariant under each
    adjoint letter of $B$. The adjoint transfer map therefore intertwines
    expansion $Z\mapsto KZK^\dagger$ with the compressed transfer map.
    Compression preserves stationarity, and cyclicity of trace preserves
    normalization. A trace-one stationary matrix for $A$ expands to one
    for $B$; uniqueness and $K^\dagger K=I$ give the claim.
\end{proof}

\begin{lemma}[Local cohomology from unique stationary bond data]
    \label{thm:spt_local_unique_stationary_class}
    \lean{MPSTensor.eventually_cohomologousTo_of_continuousOn_unique_stationary_density,
        MPSTensor.eventually_cohomologousTo_of_continuous_unique_stationary_density,
        MPSTensor.eventually_cohomologousTo_of_continuous_unital_supported_family}
    \leanok
    Let $T$ be a topological space and let $B(t)$, $t\in T$, be a
    continuous tensor family of fixed bond dimension. Let
    $\sigma(t)\geq0$ be trace-one and adjoint stationary.
    Let the minimal tensors be support compressions of $B(t)$,
    with unitary projective actions implementing a fixed exact physical
    symmetry. Suppose the compressed tensor at $t_0$ is one-site injective.
    If $\sigma$ is continuous and unique among trace-one stationary
    matrices on an open neighborhood of $t_0$, then the virtual factors
    at nearby $t$ are cohomologous to the factor at $t_0$.
    Alternatively, it suffices that $B(t)$ be unital and that its adjoint
    fixed space at $t_0$ have complex dimension one. In this case density
    continuity and local uniqueness follow automatically. Neither case
    assumes continuity of the minimal support frames or virtual actions.
\end{lemma}

\begin{proof}\leanok
    \uses{thm:spt_stationary_density_compression,
        thm:spt_unique_stationary_symmetry,
        thm:spt_local_stationary_density,
        thm:spt_local_spectral_support,
        thm:spt_supported_frame_comparison,
        thm:spt_local_invariant_compression}
    Compression preserves stationary uniqueness. Exact physical covariance
    then forces each virtual action to commute with the compressed
    density. Apply the spectral-corner argument on the open neighborhood.
    In the unital case, first construct the local unique continuous
    density and identify the given pointwise stationary choice with it.
\end{proof}

\begin{lemma}[Local continuity of a normalized stationary density]
    \label{thm:spt_local_stationary_density}
    \lean{MPSTensor.exists_local_continuous_adjointStationaryDensity_of_unital,
        MPSTensor.exists_open_continuousOn_stationaryDensity_of_unital}
    \leanok
    Let $B(t)$ be a continuous family of unital tensors of positive,
    fixed bond dimension. Suppose the adjoint transfer map at $t_0$
    has a one-dimensional fixed space over $\C$. Then there is an open
    neighborhood of $t_0$ on which it has a unique trace-one stationary
    matrix $\sigma(t)$. This matrix is positive semidefinite and continuous;
    its rank may change. Consequently any pointwise choice of trace-one
    stationary matrices agrees with this family and is continuous there.
\end{lemma}

\begin{proof}\leanok
    Write $\mathcal E_t^*$ for the adjoint transfer map and choose a
    trace-one matrix $\tau$. The linear operator
    $F_t(X)=X-\mathcal E_t^*(X)+\tr(X)\tau$ is invertible at $t_0$:
    trace preservation makes every vector in its kernel traceless, and
    the one-dimensional stationary space then forces that vector to vanish.
    Invertibility persists nearby, and $\sigma(t)=F_t^{-1}(\tau)$ is
    continuous. Its trace is one and it is stationary. Existence of a
    positive stationary state for a trace-preserving positive map
    identifies it with that state. The same inverse proves uniqueness.
\end{proof}

\begin{lemma}[Continuous normalized Perron data]
    \label{thm:spt_continuous_right_perron_data}
    \lean{MPSTensor.continuous_rightPerronDensity_of_isInjective,
        MPSTensor.exists_continuous_rightPerronPair_of_isInjective}
    \leanok
    Let $A(t)$ be a continuous family of one-site injective tensors
    in a fixed positive bond dimension. Its positive right Perron
    matrix $\rho(t)$ and positive Perron value $r(t)$ admit continuous
    choices satisfying
    \begin{align}
        \rho(t)                    & >0, \notag           \\
        \tr\rho(t)                 & =1, \notag           \\
        \mathcal E_{A(t)}(\rho(t)) & =r(t)\rho(t). \notag
    \end{align}
    In fact, any pointwise choice of positive trace-one Perron matrices
    is continuous. This is a consequence of positive-map
    Perron--Frobenius theory, applied to the given tensor family
    \cite[Theorem~6.3]{Wolf2012Quantum}.
\end{lemma}

\begin{proof}\leanok
    Perron--Frobenius theory gives a positive definite right eigenmatrix
    and a positive definite left eigenmatrix at the same positive
    eigenvalue. Pairing with the left eigenmatrix fixes the eigenvalue
    of any normalized positive right eigenmatrix. Uniqueness then fixes
    the matrix itself. Compactness of the density matrices and the
    closed eigenvector equation give continuity of the normalized choice.
    Finally $r(t)=\tr\mathcal E_{A(t)}(\rho(t))$ is continuous.
\end{proof}

\begin{lemma}[Local continuous unital normalization]
    \label{thm:spt_continuous_unital_normalization}
    \lean{MPSTensor.exists_local_continuous_unitalNormalization_of_isInjective}
    \leanok
    \uses{def:gauge_equiv}
    Let $A(t)$ be continuous in a fixed positive bond dimension and
    one-site injective at $t_0$. On an open neighborhood of $t_0$,
    there are continuous Perron data as above and a continuous tensor
    \begin{align}
        \label{eq:spt_continuous_unital_tensor}
        B^i(t)=r(t)^{-1/2}\rho(t)^{-1/2}A^i(t)\rho(t)^{1/2}.
    \end{align}
    The tensor $B(t)$ is unital and is gauge-equivalent to
    $r(t)^{-1/2}A(t)$. The continuous tensor family is a hypothesis;
    no choice from physical ground vectors is asserted here.
\end{lemma}

\begin{proof}\leanok
    \uses{thm:spt_local_invariant_compression,
        thm:spt_continuous_right_perron_data,
        thm:unital_gauge_of_pd_fixed_point}
    Injectivity persists locally. Choose the normalized Perron data
    there. Positive square roots and their inverses are continuous on
    the positive definite matrices, so~(\ref{eq:spt_continuous_unital_tensor})
    is continuous. The Perron equation gives
    $\sum_iB^i(t)B^i(t)^\dagger=I$, and the invertible matrix
    $\rho(t)^{-1/2}$ gives the stated bond conjugation.
\end{proof}

\begin{lemma}[Local continuous canonical bond data]
    \label{thm:spt_continuous_canonical_normalization}
    \lean{MPSTensor.exists_local_continuous_canonicalNormalization_of_isInjective}
    \leanok
    Under the preceding hypotheses, the neighborhood may be chosen
    so that the unital tensor $B(t)$ also has a continuous normalized
    adjoint stationary density $\sigma(t)$ satisfying
    \begin{align}
        \sigma(t)                      & \ge0, \notag       \\
        \tr\sigma(t)                   & =1, \notag         \\
        \mathcal E_{B(t)}^*(\sigma(t)) & =\sigma(t). \notag
    \end{align}
    This density is unique among all trace-one fixed matrices.
    The one-dimensional fixed-space property is derived, rather than
    assumed. These are local canonical data for the given tensor
    family; reconstruction from a physical uniform gap remains separate
    \cite{gap:spc11_spt_interpolation_upper_range}.
\end{lemma}

\begin{proof}\leanok
    \uses{thm:spt_continuous_unital_normalization,
        thm:injective_irreducible,
        thm:spt_local_stationary_density}
    Gauge equivalence and nonzero scalar normalization preserve
    injectivity. For the unital base tensor, the adjoint channel is
    trace preserving and irreducible. Its normalized stationary density
    is positive definite, and Perron--Frobenius theory makes its fixed
    space one-dimensional. Apply the local continuous stationary-density
    theorem and restrict the Perron data to the resulting smaller open
    neighborhood.
\end{proof}

% Auxiliary results for local continuous minimal realizations.
\begin{lemma}[Gram left inverses]
    \label{thm:spt_gram_left_inverse}
    \lean{Matrix.gramLeftInverse_mul, Matrix.continuous_gramLeftInverse}
    \leanok
    If a complex rectangular matrix $C$ has independent columns, then
    $H=(C^\dagger C)^{-1}C^\dagger$ satisfies $HC=I$.
    For a continuous family with independent columns, $H$ is continuous.
\end{lemma}
\begin{proof}\leanok
    The Gram matrix is positive definite and hence invertible.
    Matrix inversion is continuous on the invertible matrices.
\end{proof}

\begin{lemma}[Coordinates of a bilinear map]
    \label{thm:spt_bilinear_coordinates}
    \lean{Matrix.toLinearMap₂'_apply_mulVec_prod}
    \leanok
    Let $M$ have columns indexed by pairs $(a,b)$. The associated bilinear
    map satisfies $\mu(x,y)=M(x\otimes y)$, where
    $(x\otimes y)_{(a,b)}=x_a y_b$.
\end{lemma}
\begin{proof}\leanok
    Expand both sides in the coordinate basis.
\end{proof}

\begin{lemma}[Polynomial correction of an idempotent]
    \label{thm:spt_newton_idempotent}
    \lean{idempotent_newton_defect, idempotent_newton_sub,
        idempotent_newton_twice_sub_one, idempotent_newton_increment_bound,
        idempotent_newton_tendsto_of_small_defect}
    \leanok
    In a complete normed ring with $\|1\|\le1$, set
    $x_{n+1}=3x_n^2-2x_n^3$ and $d_n=x_n^2-x_n$.
    Then $d_{n+1}=d_n^2(4d_n-3)$ and
    $x_{n+1}-x_n=-(2x_n-1)d_n$.
    If $\|d_0\|\le1/16$, the sequence converges to an idempotent.
    More precisely,
    \begin{align}
        \|x_{n+1}-x_n\| & \le
        \|2x_0-1\|\,\|d_0\|\,2^{-n}.
        \notag
    \end{align}
\end{lemma}
\begin{proof}\leanok
    Polynomial expansion gives the identities, including
    $2x_{n+1}-1=(2x_n-1)(1-2d_n)$.
    Induction gives $\|d_n\|\le4^{-n}\|d_0\|$ and
    $\|2x_n-1\|\le2^n\|2x_0-1\|$. The increments are summable.
    Completeness gives the limit; continuity of multiplication makes it
    idempotent.
\end{proof}

\begin{lemma}[Continuous local idempotents]
    \label{thm:spt_continuous_nearby_idempotent}
    \lean{Matrix.exists_local_continuous_idempotent_in_adjoin}
    \leanok
    Let $X(t)$ be a continuous family of complex square matrices,
    with $X(t_0)$ idempotent. On an open neighborhood of $t_0$ there is
    a continuous family of idempotents $P(t)$, with $P(t_0)=X(t_0)$,
    belonging pointwise to the unital algebra generated by $X(t)$.
    No Hermiticity assumption is required.
\end{lemma}
\begin{proof}\leanok
    \uses{thm:spt_newton_idempotent}
    Shrink the neighborhood so that the idempotency defect is at most
    $1/16$ and $2X-I$ is bounded. Apply the preceding correction in the
    complete normed algebra of bounded continuous matrix functions.
    Each iterate belongs to the generated algebra, which is closed
    by finite dimensionality. At $t_0$ every iterate equals $X(t_0)$.
\end{proof}

\begin{lemma}[Rank of an idempotent]
    \label{thm:spt_idempotent_rank}
    \lean{Matrix.trace_eq_rank_of_idempotent,
        Matrix.exists_open_constantRank_of_idempotent}
    \leanok
    An idempotent complex matrix has trace equal to its rank.
    Consequently the rank of a continuous family of idempotents is
    locally constant, including for non-Hermitian idempotents.
\end{lemma}
\begin{proof}\leanok
    The associated linear map is the projection onto its range.
    Its trace is the dimension of that range. Continuity of the trace
    and discreteness of the integer ranks give local constancy.
\end{proof}

\begin{lemma}[Continuous frames for a matrix range]
    \label{thm:spt_continuous_idempotent_frame}
    \lean{Matrix.exists_open_continuous_rangeFrame_of_constantRank,
        Matrix.exists_local_continuous_rangeFrame_of_idempotent}
    \leanok
    A continuous matrix family of constant rank $D$ admits a local
    continuous frame $C(t)=P(t)F$ for its range, with a fixed matrix $F$.
    The matrices $H(t)=(C(t)^\dagger C(t))^{-1}C(t)^\dagger$ are continuous
    and satisfy $H(t)C(t)=I_D$.
    For an idempotent family it suffices that $P(t_0)$ have rank $D$;
    in addition $P(t)C(t)=C(t)$.
\end{lemma}
\begin{proof}\leanok
    \uses{thm:spt_trace_quotient_fixed_section, thm:spt_gram_left_inverse,
        thm:spt_idempotent_rank}
    Choose a section at the base parameter. Injectivity persists nearby.
    Equality of dimensions identifies its range with the matrix range.
    Apply the Gram left inverse and the preceding local rank result.
\end{proof}

\begin{lemma}[Dimension of a primitive left ideal]
    \label{thm:spt_primitive_left_ideal_dimension}
    \lean{Matrix.finrank_range_mulRight_single,
        Matrix.finrank_range_rightMul_of_matrixAlgebra}
    \leanok
    In $\MN{D}$ the left ideal $\MN{D}E_{ii}$ has
    dimension $D$. The same holds under any algebra
    isomorphism.
\end{lemma}
\begin{proof}\leanok
    The ideal consists of matrices supported in column $i$.
    That column is an arbitrary vector in $\C^D$.
\end{proof}

\begin{lemma}[Continuous primitive left ideals of a trace quotient]
    \label{thm:spt_continuous_trace_quotient_primitive}
    \lean{Matrix.exists_local_continuous_traceQuotientPrimitive}
    \leanok
    Suppose a continuous ambient tensor family has two-site and three-site
    traces proportional, with nonzero factors, to those of pointwise
    injective tensors of constant positive bond dimension $D$.
    Its locally reconstructed trace-quotient algebra admits a continuous
    idempotent $p(t)$ whose right multiplication has rank $D$.
    Both the multiplication coordinates and this right-multiplication
    matrix are continuous. No continuous choice of the pointwise minimal
    tensors or algebra identifications is assumed.
\end{lemma}
\begin{proof}\leanok
    \uses{thm:spt_trace_quotient_continuous_unit,
        thm:spt_continuous_nearby_idempotent,
        thm:spt_primitive_left_ideal_dimension, thm:spt_idempotent_rank}
    In the base matrix algebra choose a diagonal matrix unit and
    pass it to the quotient. Keep its coordinates fixed and
    correct its left-regular action to a nearby idempotent.
    The generated-algebra property places the correction in the
    left-regular representation. Apply it to the continuously chosen
    unit to obtain $p(t)$. The principal left ideal has rank $D$ at
    the base parameter, hence nearby by local constancy.
\end{proof}

\begin{definition}[Action in a frame of a left ideal]
    \label{def:spt_left_ideal_frame_action}
    \lean{Matrix.leftIdealFrameAction}
    \leanok
    For a bilinear product $\mu$, a frame $J$ of a principal left ideal,
    and a left inverse $H$, define $L(x)v=H\mu(x,Jv)$.
\end{definition}

\begin{lemma}[Matrix realization in a left-ideal frame]
    \label{thm:spt_left_ideal_frame_realization}
    \lean{Matrix.leftIdealFrameAction_mulVec,
        Matrix.leftIdealFrameAction_intertwines,
        Matrix.exists_inner_of_framed_leftAction,
        Matrix.exists_inner_of_leftIdealFrame,
        Matrix.exists_linearEquiv_leftIdealFrameAction,
        MPSTensor.gaugeEquiv_leftIdealFrameAction}
    \leanok
    \uses{def:spt_left_ideal_frame_action}
    Suppose the quotient algebra is isomorphic to $\MN{D}$,
    with $D>0$, and $J$ identifies $\C^D$ with a principal left
    ideal, with $HJ=I$. The framed action is an algebra isomorphism
    onto $\MN{D}$ and differs from the given identification
    by an inner conjugation. It preserves traces. Applying it to
    a family of letters therefore preserves that family's bond gauge
    equivalence.
\end{lemma}
\begin{proof}\leanok
    Associativity makes the left ideal invariant and gives
    $JL(x)v=\mu(x,Jv)$. This action is multiplicative and unital.
    Composing with the inverse of the given algebra identification
    gives a unital endomorphism of the full matrix algebra.
    The Skolem--Noether theorem makes it an inner conjugation.
\end{proof}

\begin{definition}[Coordinates of reconstructed letters]
    \label{def:spt_trace_quotient_letter_coordinates}
    \lean{Matrix.traceQuotientLetterCoordinates}
    \leanok
    For the two-site trace form $G$ and an injective section $F$, set
    $q_i=((GF)^\dagger GF)^{-1}(GF)^\dagger G e_i$.
\end{definition}

\begin{lemma}[Recovery of the letters]
    \label{thm:spt_trace_quotient_letter_recovery}
    \lean{Matrix.traceQuotientLetterCoordinates_recover,
        MPSTensor.gaugeEquiv_traceQuotientLetters,
        Matrix.traceQuotientLetters_of_normalized_linearEquiv}
    \leanok
    \uses{def:spt_trace_quotient_letter_coordinates}
    If $G_{ji}=\alpha_2\tr(A^iA^j)$ and the normalized
    quotient identification satisfies
    $E(x)=\kappa\sum_i(Fx)_i A^i$, with $\alpha_2\kappa\ne0$,
    then $E(q_i)=\kappa A^i$. In any left-ideal frame the represented
    letters are gauge-equivalent to $\kappa A$.
\end{lemma}
\begin{proof}\leanok
    \uses{thm:spt_gram_left_inverse, thm:spt_left_ideal_frame_realization,
        thm:spt_trace_quotient_contraction, thm:trace_pairing_injective}
    Injectivity of the minimal tensor makes its trace pairing
    nondegenerate. The Gram left inverse recovers the unique letter
    coordinates. Apply the preceding inner-conjugacy statement.
\end{proof}

\begin{lemma}[Continuity of represented letters]
    \label{thm:spt_continuous_left_ideal_letters}
    \lean{Matrix.continuousOn_leftIdealFrameAction,
        Matrix.continuousOn_traceQuotientLetterCoordinates}
    \leanok
    \uses{def:spt_left_ideal_frame_action,
        def:spt_trace_quotient_letter_coordinates}
    Continuous multiplication coefficients, frames, left inverses,
    and letter coordinates give continuous represented letters.
    The letter coordinates are continuous wherever the fixed section
    is injective.
\end{lemma}
\begin{proof}\leanok
    \uses{thm:spt_gram_left_inverse}
    The formulas are finite sums and matrix products. The only inverse
    is the Gram inverse, which is continuous on this neighborhood.
\end{proof}

\begin{lemma}[Finite traces from positive-length periodic rays]
    \label{thm:spt_positive_ray_finite_traces}
    \lean{MPSTensor.exists_nonzero_two_three_trace_scalars_of_samePositiveMpvRay}
    \leanok
    If an ambient tensor and a one-site injective tensor of positive bond
    dimension have the same periodic rays at every positive length,
    their two-site and three-site word traces are proportional by
    nonzero scalars $\alpha_2$ and $\alpha_3$.
\end{lemma}
\begin{proof}\leanok
    \uses{thm:mpoalg_length_dependent_phase_power}
    Specialize the ray equality to lengths two and three and evaluate
    each vector in the physical coordinate basis.
\end{proof}

\begin{lemma}[Local continuous minimal realization from finite traces]
    \label{thm:spt_continuous_minimal_finite_trace_realization}
    \lean{MPSTensor.exists_local_continuous_minimalTensor_of_two_three_trace_eq}
    \leanok
    Let $B(t)$ be continuous in a fixed ambient bond dimension.
    Suppose pointwise one-site injective tensors $A(t)$ have a fixed
    positive bond dimension $D$ and satisfy
    \begin{align}
        \tr(B^i(t)B^j(t))       & =\alpha_2(t)\tr(A^i(t)A^j(t)),
        \notag                                                   \\
        \tr(B^i(t)B^k(t)B^j(t)) & =
        \alpha_3(t)\tr(A^i(t)A^k(t)A^j(t)),
        \notag
    \end{align}
    with nonzero factors. Around each parameter there is a continuous
    one-site injective tensor $C(t)$ of dimension $D$ that is
    gauge-equivalent to $(\alpha_3(t)/\alpha_2(t))A(t)$.
    The pointwise tensors, scalar factors, and gauges need not be
    continuous. This is a reconstruction at constant minimal
    dimension; the two trace identities alone make no assertion about
    the ambient tensor's longer periodic rays.
\end{lemma}
\begin{proof}\leanok
    \uses{thm:spt_continuous_trace_quotient_primitive,
        thm:spt_continuous_idempotent_frame,
        thm:spt_left_ideal_frame_realization,
        thm:spt_trace_quotient_letter_recovery,
        thm:spt_continuous_left_ideal_letters}
    Recover the continuous quotient product and unit, choose a
    continuous primitive left ideal, and construct a continuous frame
    for that ideal. Left multiplication in the frame represents
    the recovered letters continuously. The pointwise inner-conjugacy
    theorem identifies them with the normalized minimal tensor.
\end{proof}

\begin{lemma}[Local continuous minimal realization of periodic rays]
    \label{thm:spt_continuous_minimal_ray_realization}
    \lean{MPSTensor.exists_local_continuous_minimalTensor_of_samePositiveMpvRay}
    \leanok
    Suppose the continuous ambient family $B(t)$ has, at every parameter,
    the same positive-length periodic rays as a pointwise one-site
    injective tensor $A(t)$ of a fixed positive bond dimension.
    Locally those same rays admit continuous one-site injective
    representatives of that bond dimension. No continuity of the
    given pointwise representatives is assumed.
\end{lemma}
\begin{proof}\leanok
    \uses{thm:spt_positive_ray_finite_traces,
        thm:spt_continuous_minimal_finite_trace_realization,
        thm:spt_exact_phase_gauge_invariance}
    Obtain nonzero two-site and three-site trace factors from the
    full ray equality. Apply the preceding reconstruction.
    Nonzero scalar multiplication and bond conjugation preserve
    every positive-length periodic ray.
\end{proof}

\begin{lemma}[Local continuous canonical realization of periodic rays]
    \label{thm:spt_continuous_canonical_ray_realization}
    \lean{MPSTensor.exists_local_continuous_unitalTensor_of_samePositiveMpvRay}
    \leanok
    Under the preceding hypotheses, the local representatives can
    be chosen continuous, one-site injective, and unital, with a
    continuous normalized adjoint stationary density $\sigma(t)$.
    This density is positive semidefinite, has trace one, and is unique
    among all trace-one fixed matrices of the adjoint transfer map.
    The minimal bond dimension is constant and positive. Reconstruction
    through a change of minimal dimension remains separate
    \cite{gap:spc11_spt_interpolation_upper_range}.
\end{lemma}
\begin{proof}\leanok
    \uses{thm:spt_continuous_minimal_ray_realization,
        thm:spt_continuous_canonical_normalization,
        thm:spt_exact_phase_gauge_invariance}
    Apply continuous Perron normalization to the reconstructed family
    on a smaller neighborhood. Its scalar gauge preserves the full
    ray family. Restrict the continuous unital tensor and stationary
    density to the resulting ambient open neighborhood.
\end{proof}
